% !TEX program = xelatex
% ============================================================
%  Philosophy of Intimacy and the Theory of Justice · Paper IX
%  The Sustainability of Intimacy:
%  Poetic Generativity and Generative Justice
%
%  Uses the shared series style (serendip-paper.sty). XeLaTeX +
%  Noto Serif CJK SC for occasional Chinese terms.
% ============================================================
\documentclass[11pt,a4paper]{article}
\PassOptionsToPackage{table}{xcolor}

% ---------- Fonts (CJK fallback for occasional Chinese terms) ----------
\usepackage{fontspec}
\newfontfamily\cjk{Noto Serif CJK SC}
\newcommand{\zh}[1]{{\cjk #1}}
\XeTeXlinebreaklocale "zh"
\XeTeXlinebreakskip = 0pt plus 0.1pt

% natbib (numbered) before serendip-paper (hyperref loads inside it)
\usepackage[backend=biber, style=authoryear, sorting=nyt, maxcitenames=2, uniquelist=false]{biblatex}
\addbibresource{../../bib/master.bib}
\usepackage{serendip-paper}

% colour used on the title page (headings/links come from serendip-paper)
\definecolor{inksoft}{HTML}{342E37}

% ---------- Math ----------
\newcommand{\holo}{\operatorname{hol}}

% ---------- Tables (beyond the shared style) ----------
\usepackage{tabularx}
\usepackage{pdflscape}
\usepackage{ragged2e}
\renewcommand{\arraystretch}{1.3}
\newcolumntype{Y}{>{\RaggedRight\arraybackslash}p{3.1cm}}

% ---------- Paper-specific emphasis & helpers ----------
\newcommand{\emc}[1]{\textcolor{rose}{\oldtextit{#1}}}
\newcommand{\emb}[1]{\textcolor{rose}{\oldtextbf{#1}}}
\newcommand{\goldrule}{\noindent\textcolor{gold}{\rule{\linewidth}{0.6pt}}\par\vspace{0.4em}}
\newcommand{\todo}[1]{\textcolor{gray}{\small[TODO: #1]}}
\newcommand{\xref}[1]{\textcolor{rose}{\S\ref{#1}}}
\newcommand{\pinkref}[1]{\textcolor{rose}{#1}}

% A claim set off from the prose: a thin gold rule above and below,
% used sparingly for the load-bearing propositions of the paper.
\newenvironment{claim}[1][]{%
  \par\vspace{0.6em}\noindent\textcolor{gold}{\rule{\linewidth}{0.4pt}}\par\vspace{0.4em}
  \ifx\relax#1\relax\else\noindent\emb{#1.}\hspace{0.4em}\fi\ignorespaces}{%
  \par\vspace{0.4em}\noindent\textcolor{gold}{\rule{\linewidth}{0.4pt}}\par\vspace{0.6em}}

% Inline prose case: a run-in italic label.
\newenvironment{casebox}[1]{%
  \par\vspace{0.4em}\noindent\textit{\textcolor{rose}{An illustrative case --- #1.}}\hspace{0.4em}\ignorespaces}{%
  \par\vspace{0.4em}}

% ============================================================
\begin{document}

% ---------- Title page ----------
\thispagestyle{empty}
\begin{center}
\vspace*{2cm}
{\footnotesize\color{rosedeep}\scshape Philosophy of Intimacy and the Theory of Justice \quad$\cdot$\quad Paper IX}\\[0.3cm]
\coverrule\\[1.2cm]
{\LARGE\bfseries\color{warnred} The Sustainability of Intimacy}\\[0.7cm]
{\large\color{warnred} Poetic Generativity and Generative Justice}\\[0.8cm]
{\small\color{gold} [\,Working Draft\,]}\\[0.8cm]
{\normalsize\color{inksoft} On the Good and the Vicious Circle, the Geometry of Sublimation, and a Eudaimonics of the Relational Field}\\[2cm]
\coverrule[0.25]\\[0.6cm]
{\normalsize\color{inksoft} Wanhong}\\
\vfill
{\footnotesize\color{gold} Working draft --- not for citation or circulation}
\coverfootnote
\end{center}
\newpage

% ---------- Epigraph ----------
\thispagestyle{empty}
\vspace*{2.4cm}
\begin{center}
\begin{minipage}{0.74\linewidth}
\centering\itshape\color{rose}
{\cjk 但愿人长久，千里共婵娟。}\\[0.3em]
So let us wish each other long life, and share, though a thousand\\miles apart, the same fair moon.\\[0.8em]
{\cjk 夫物芸芸，各复归其根。}\\[0.3em]
The myriad things in their teeming each return to their root.\\[0.8em]
{\cjk 生而不有，为而不恃，长而不宰，是谓玄德。}\\[0.3em]
To generate without possessing, to act without presuming,\\to foster without ruling: this is called the dark virtue.\\[0.6em]
{\upshape\footnotesize\color{gold} ---{\cjk 苏轼}\ Su Shi\,;\ {\cjk《道德经》}\ \emph{Daodejing}, ch.~16 \& 51}
\end{minipage}
\end{center}
\newpage

% ---------- Dedication ----------
\thispagestyle{empty}
\vspace*{4.5cm}
\begin{center}
\begin{minipage}{0.7\linewidth}
\centering\itshape\color{rose}
For her---\\[0.7em]
\upshape the girl of the forest,\\[0.3em]
who loves the woods and the long road of travelling;\\[0.7em]
\itshape pure, and kind, and touched with the divine.\\[1.0em]
{\upshape\footnotesize\color{gold} To the one I love most.}
\end{minipage}
\end{center}
\newpage

% ---------- Abstract ----------
\thispagestyle{empty}
\noindent\textbf{\large\color{rosedeep} Abstract}\\[0.6em]
\noindent
The earlier papers of this series asked how a relation \emph{begins}; this one asks how it does not \emph{end}---under what conditions a generative bond reproduces itself across time. Its problem is that self-continuation, taken alone, is morally indifferent: capital, the rose, and the parasitic bond all continue themselves no less than love does. Eglash's recursive return of value is therefore necessary but not sufficient. The paper supplies the missing temporal dimension through an axiom of sustainability cast in the language of \emph{geometric phase}---the vicious cycle a circle of zero holonomy, the good cycle a spiral that returns to its root while sublimating---and then asks how, through signs, such a spiral might be fostered. Its constructive thesis is that the \emph{appropriate poem}---the non-possessive interpretive cycle that generates real value---is the form of the good circle, and that generative justice, in the intimate domain, is for this reason necessarily poetic.


\vspace{1.2em}
\noindent{\color{rosedeep}\textbf{Keywords:}}\, \textcolor{inksoft}{sustainability of intimacy; generative justice; geometric phase; the good and the vicious circle; non-possession; poetic generativity.}
\newpage

\tableofcontents
\newpage

% ============================================================
%  Body
% ============================================================
% ===== §1 Prelude: the rose in the garden =====
\section{Prelude: The Rose in the Garden}
\label{p09:sec:prelude}

\noindent
In a garden, someone said a small thing that became the seed of this paper. The rose, they remarked, is beautiful; and yet, however hard it labours, it can hold its bloom for only a week or so.

It is the kind of observation one lets pass---a pleasantry about the brevity of beautiful things, the sort of melancholy that attaches to gardens and is half the reason one keeps them. But it lodged, and would not leave, because it names---without meaning to---the precise difficulty this paper is about. The remark seems to set beauty against duration, as though the rose were beautiful \emph{despite} being brief, the brevity a flaw in an otherwise perfect thing. The truth is harder and more interesting: the beauty and the brevity are not two facts but one. \emb{The rose blooms so resplendently \emph{because} it pours the whole of its resource into a single, unrepeatable flowering, holding nothing back against the morrow.} Its resplendence is the very form of its unsustainability. To wish the bloom both as intense and as lasting is to wish a contradiction---to want the flower to spend itself utterly and to keep itself in reserve at once. There is, here, a tragic trade-off that the paper will meet again and again, and that it had better not pretend to abolish: intensity and persistence pull, very often, against one another, and the most intense things are, for that reason, the briefest.

And the same structure, I want to claim, is at work in domains that share no soil with the rose. The intimate bond, at its most incandescent, has the rose's economy: the passion that consumes the whole of itself in its flowering is, precisely thereby, the passion that cannot last, and the lover who demands that the first ardour both blaze and endure demands the rose's contradiction. The logic of capital has it too: the accumulation that drives toward an ever-greater intensity of valorization is the accumulation that exhausts the conditions---the worker, the soil, the future---on which its own continuation depends, blazing brightest as it nears the exhaustion of what feeds it. And the dynamics of the cosmos, at the largest scale, show the same face: the structures that burn most brightly are the most prodigal, and what endures is rarely what shines hardest. Across all three---and the recurrence is itself a reason to attend---one finds a generativity that takes its own continuation as its end, production pointing back at production itself (\zh{反者道之动}, \emph{reversal is the movement of the Dao}), and yet a generativity that may, like the rose, burn itself out in the very intensity by which it lives.

\emc{So the question of this paper is not the one the remark seems to pose.} It is not \emph{how to keep the flower from wilting}---that wish, as we shall see, is the very thing that rots the root. It is the question hidden beneath the remark, and harder to hear: \emph{when the flower has wilted, how does the root go on?} For the observation, true of the flower, is false of the rosebush. The single bloom is mortal and exhaustive, spending all in a week; but the bush flowers year upon year, and it endures not in spite of the flower's death but \emph{by} it---the fallen petals returning their nutrient to the root, this year's wilting the condition of next year's bloom. Here is the distinction on which the whole paper turns, and it is worth stating in its starkest form at the outset:

\begin{claim}[Sustainability is the nesting of finitude, not its abolition]
Sustainability is not the abolition of finitude---not a flower that never falls, which is the bad infinite, the fantasy of perpetual growth that the paper will identify with the vicious circle. Sustainability is the \emph{nesting} of finitude within a higher cycle of regeneration: the flower must fall \emph{so that} the root may go on. A bloom held past its hour does not extend the rose's life; it is the refusal of the fall that kills the bush. The good circle, far from excluding death, requires it---requires that the flower be spent and returned---and what we will come to call the vicious circle is, very often, the one that cannot let its flower fall.
\end{claim}

This reframing carries a consequence the paper will draw out at length, and which it is worth stating now in its sharpest form, for it overturns an assumption so ordinary we scarcely know we hold it. We assume that the threat to a beautiful thing is neglect---that what is precious withers for want of tending, and that more care is therefore always the safer error. The rose teaches otherwise. \emb{The deepest unsustainability comes not from too little tending but, very often, from a tending driven by the fear of unsustainability itself.} The one who, dreading that the rose will wilt, waters it without cease, changes its water, fusses at its roots---that one does not save the flower; he drowns the root, and so manufactures the very withering he feared. This is no incidental cruelty of fate but a structure, and the paper will meet it in every domain it touches: in capital's dread of the end of growth, which drives the extraction that exhausts its base; in the lover's dread of loss, which drives the control and the clutching that rot the bond; in the family's dread of dissolution, which hardens into the very rigidity that breaks it. \zh{为者败之，执者失之}---``the one who acts upon it ruins it, the one who grasps it loses it.'' The fear of an end, made into action, produces the end it fears.

There is, then, a quiet reversal of the imagination at work in this paper from its first page. The one who said the rose blooms but a week was speaking, lovingly, of the flower, and of the sadness of its falling. This paper means to speak of the root---of what is not seen, of what the falling feeds, of the slow regeneration that the flower's brevity makes possible and the anxious hand makes impossible. It is, in a sense, a paper about learning to let the flower fall: not from indifference to its beauty, but from a deeper attention to the root that the beauty, rightly let go, sustains.

\emc{A word, finally, on where this paper stands in the series.} The first eight papers asked, in the main, how a relation \emph{begins}: the mother paper established the ontology of relational being, that the subject is generated in and through relation rather than entering it already complete---a lineage that runs back through Buber's account of the I--Thou as the primary word that constitutes both terms \parencite{buber1970}; \pinkref{Paper~V} descended to the joint attention by which two consciousnesses first constitute one another; \pinkref{Paper~VI} cut into the single decisive moment of the proposal and exposed the epistemic injustice latent within it; \pinkref{Paper~VII} took the dyad's first turning toward the world, the diplomacy by which it faces the external relational field; \pinkref{Paper~VIII} treated the lyric's address to the big Other. Each took, as its object, a \emph{founding moment}---an opening, an inception, a first constitution. \emb{The present paper takes, for the first time, the opposite object: not how a relation begins, but how it does not end.} It turns from inception to maintenance, from the founding act to the long reproduction that must follow it, from the question \emph{how to begin} to the question \emph{how not to end}---which is, we shall find, the question of the root rather than the flower, and the first time the series has asked it. This is its place, and its novelty, in the programme of Generative Relational Being.
      % §1 Prelude: the rose in the garden
% ===== §2 Thesis: generativity takes its own continuation as its end =====
\section{Thesis: Generativity Takes Its Own Continuation as Its End}
\label{p09:sec:thesis}

\noindent
The paper begins from a metaphysical conjecture, and states it at once, baldly, so that the antithesis may have something firm to push against. \emb{Generativity takes its own continuation as its end. Production points back at production itself.} The aim of a generative process is not some terminal product at which it would come to rest, but the perpetuation of generating---the cycle's end is the cycle. A generativity oriented toward a final product would, upon attaining it, stop; the generativity this paper is about does not aim at a state it could reach and rest in, but at its own ongoingness. Its telos is itself.

This is a strong claim, and its strangeness should not be smoothed over. We are accustomed to think of processes as means to ends---production as the making of a product, desire as the seeking of an object, labour as the securing of a wage---and on that picture the process is subordinate, justified by the end it serves, and complete when the end is attained. The conjecture inverts this: in the generativity that concerns us, the ``product'' is not the point at which the process is justified and stops, but the occasion for the process to continue. The product is for the sake of the producing, not the producing for the sake of the product. This is what is meant by calling such generativity \emph{autotelic}: it carries its end within itself, as its own continuation, and is for that reason not completed by any external attainment.

The conjecture is not arbitrary, and its warrant is the constancy with which it recurs across domains that share no subject matter. Four such recurrences are worth setting out, for the breadth of the pattern is the chief reason to take it for something like a metaphysical ground rather than a local fact about one or another region of the world.

\emc{In the dynamics of the cosmos}, the structures that persist are, overwhelmingly, the self-reproducing ones. A dissipative structure---a flame, a whirlpool, a living cell---maintains itself far from thermodynamic equilibrium not by reaching some final stable state but by continually regenerating the very gradients and flows that constitute it; it persists by producing the conditions of its own persistence. The autocatalytic cycle, the chemical paradigm of life's origin, is a set of reactions whose product is the catalyst of its own production---a loop whose output is the cause of its output. In both, what endures is precisely what makes more of its own enduring; the structures that aim at no terminal state, but only at their own continuation, are the ones the universe keeps. The autotelic is not an oddity of human striving but a feature of what persists at all.

\emc{In the dialectic of desire}, as psychoanalysis describes it, the same structure governs the subject. Desire does not aim at a final satisfaction that would extinguish it; were it ever fully satisfied, it would cease, and the subject of desire would dissolve. Beneath every particular object---this person, this attainment, this possession---desire aims at its own continuation, which is why the object attained is, with a regularity that is structural and not accidental, found wanting, and desire moves on. This is not the failure of desire but its success: desire is organized to defer the satisfaction that would end it, to keep itself going by never quite arriving. Lacan's distinction of desire from demand and need turns on just this: need has an object that satisfies and stops it, but desire, beneath demand, aims past every object at the perpetuation of desiring itself. The subject lives by a wanting whose true end is to go on wanting.

\emc{In the structure of dialectic itself}---thesis, antithesis, synthesis---the autotelic form appears as the engine of the movement. The synthesis that resolves a contradiction is not a resting place but the thesis of a further movement; the \emph{Aufhebung} that lifts a tension into a higher unity does not end the dialectic but generates its next turn. A dialectic that reached a final synthesis and stopped would not be a dialectic but a death; the living dialectic is the one whose every resolution opens the next contradiction, whose end is not a final synthesis but the continuation of synthesizing. Here, in the most abstract register, is the same structure the rose and the cell display: a process whose product is the occasion of its own continuation.

\emc{And in Eglash's generative justice}, this conjecture finds the political-economic and ethical articulation on which the paper will lean throughout. Eglash's insight is that value, in an unalienated economy, does not flow outward to be extracted and accumulated at a distant centre, but \emph{circulates back to those who generate it}, to be reinvested in the conditions of further generation \parencite{eglash2016}. The recursive, self-similar structure on which this regeneration leans is the through-line of Eglash's earlier study of indigenous fractal design \parencite{eglash1999}. Where the dominant economies---he argues that both the capitalist and the state-socialist are, in this, alike---are organized around the \emph{extraction} of value from its producers to a centre, the generative alternative keeps value circulating at the grassroots, among and back to those whose labour, care, and creativity produce it. Generativity, here, perpetuates itself by recursion: the value it produces returns to nourish the producers, who produce again, in a loop that is its own end. \emb{This is the spatial criterion of generative justice: that value, unalienated, return to its creators.} And the recursive structure of regeneration is, in Eglash's analysis, all but universal---generativity, to perpetuate itself at all, must lean on a recursive structure that returns its product to the conditions of its own renewal.

Drawing the four together:

\begin{claim}[The metaphysical ground, as point of departure]
Generativity takes its own continuation as its end---an autotelic, self-referential structure that recurs across the dynamics of the cosmos (the self-reproducing dissipative structure), the dialectic of desire (the wanting whose end is to go on wanting), the structure of dialectic itself (the synthesis that is the next thesis), and the political economy of generative justice (the recursive return of value to its creators). The constancy of the recurrence warrants treating it as something like a metaphysical ground. This is the paper's \emph{thesis}, its point of departure.
\end{claim}

But it must be stressed that the thesis is offered as ground and not as conclusion, for the temptation here is real and the whole argument of the paper turns on refusing it. To establish that generativity is self-perpetuating, and that justice consists in the recursive return of value to its creators, can seem already to have delivered an ethics: let value circulate back to those who make it, refuse the extraction that carries it off to a centre, and all is well. The labour of this paper is to show that this is not enough---that the recursive return of value is necessary to the good circle but does not suffice for it, because the bare fact of self-continuation, even self-continuation \emph{with} return, cannot by itself tell the good circle from the vicious one. For the most rapacious cycles are self-continuing too, and some of them return value to a circle of creators while exhausting everything outside it. To that insufficiency---the moral indifference of mere self-continuation---the antithesis now turns.
       % §2 Thesis: generativity takes its own continuation as its end
% ===== §3 Antithesis: self-continuation is morally indifferent =====
\section{Antithesis: Self-Continuation Is Morally Indifferent}
\label{sec:antithesis}

\noindent
The thesis stands; and against it the antithesis raises a single, stubborn law. \emb{Self-continuation, in itself, is morally indifferent.} That a cycle perpetuates itself tells us nothing of its goodness, for the most rapacious cycles perpetuate themselves no less than the most generous---and some of them perpetuate themselves precisely \emph{by} their rapacity. The autotelic structure that the thesis identified as a metaphysical ground is, taken alone, no mark of the good at all; it is, if anything, the formal signature of the very thing the paper means to criticize. This is the hard turn the argument must take before it can build: the ground it has just laid does not yet bear any ethical weight, and to pretend otherwise would be to mistake a structure for a value.

Consider the cases, set deliberately side by side, for their likeness is the whole of the point.

\emc{Capital} is the purest case, and the one in which the autotelic structure is most exactly visible. Marx's formula for capital is not C--M--C---commodity exchanged, by way of money, for another commodity, to satisfy a need that ends the movement---but M--C--M$'$: money advanced to become more money, a movement whose end is not any commodity, not the satisfaction of any need, but its own continuation at an enlarged scale \citep{marx1976}. Capital, in Marx's striking phrase, is value in process, value that has become the ``automatic subject'' of its own movement, and its aim is nothing but its own valorization---accumulation for the sake of accumulation, production for the sake of more production. This is self-continuation in its most exact and most relentless form: production pointing at production, the cycle whose only end is the cycle. And it is precisely the thing the paper means to criticize. Here is the first and sharpest blow to any easy reading of the thesis: the autotelic form, far from being the mark of the good, is the formal signature of what the paper will come to call the vicious circle. M--C--M$'$ satisfies the thesis to the letter, and it is the model of extraction.

\emc{The rose} is the case carried over from the prelude, and it shows that the indifference of self-continuation is not peculiar to the economic. The rose's blooming is self-continuation of a kind---the plant's drive to flower, to set seed, to flower again, a generativity oriented toward its own perpetuation through the generations. And yet, at the scale of the single flower, it is exhaustive: the resplendence is bought by a total expenditure that holds nothing in reserve, and the bloom that spends everything is, for that reason, the bloom that cannot last. The rose perpetuates the species by spending the flower; and whether this spending is good or ill---whether it is the generous fall that feeds the root or the anxious blaze that exhausts it---the bare fact of perpetuation cannot say. Self-continuation here wears neither a virtuous nor a vicious face; it is, in itself, indifferent between them.

\emc{The parasitic bond} is the case nearest the paper's true subject, and the one that should most disturb. A relation in which one party's self-continuation is purchased by the steady consumption of the other---the bond that endures, sometimes for decades, precisely by extracting from one partner the resource that sustains the other---perpetuates itself as surely as any love, and often more stably, for the extracted party's depletion can be made to look like devotion and the arrangement thereby stabilized. It is a cycle; it does not stop; by the thesis's criterion it is a fully successful generativity. And it is vicious. That a bond lasts, that it reproduces itself across years, that value circulates within it---none of this distinguishes it from love, and the failure to see this is the failure the antithesis exists to correct.

The three are formally alike: each is a self-continuing cycle, each satisfies the thesis, and the likeness is exactly what exposes the insufficiency of the thesis---and, with it, the insufficiency of Eglash's spatial criterion taken alone.

\begin{claim}[The insufficiency of the spatial criterion]
The recursive return of value---Eglash's spatial criterion---is necessary to the good circle but does not suffice for it. For it cannot, by itself, exclude two failures. It cannot exclude the \emph{fair stasis}: a cycle in which value duly returns to its creators, and which is in this sense just, yet which merely repeats, generating no surplus, climbing nowhere---the bond grown fair and lifeless, conserved but not sublimating, a circle and not a spiral. And it cannot exclude the \emph{counterfeit closure}: a cycle that appears to return value to its creators while in fact sustaining itself by extraction from an outside it has hidden from view. Self-continuation, even self-continuation \emph{with} return, is not yet the good.
\end{claim}

The two failures deserve to be drawn out, for the rest of the paper is, in large part, the apparatus for detecting them.

\emc{The fair stasis} is the gentler failure, and the one the spatial criterion is least equipped to see, since by the spatial criterion it is not a failure at all. Imagine a bond in which value genuinely returns to both parties, neither extracts from the other, the distribution is scrupulously just---and nothing grows. The cycle turns and returns to exactly where it began; each round reproduces the last without remainder; the relation is fair, and dead. This is the circle of zero holonomy that the synthesis will name: a return that returns to the same, accruing nothing, climbing nowhere. The spatial criterion pronounces it just and falls silent, for it has no temporal dimension with which to register that a just cycle may yet be a stagnant one. That a relation can be fair and lifeless at once is the first thing the spatial criterion cannot say---and the first reason a temporal criterion is needed.

\emc{The counterfeit closure} is the graver failure, and it deserves the closer look, for it is the structural heart of the antithesis and the thing the whole paper must learn to detect. Here is the puzzle it poses: if the vicious cycle is truly unsustainable, how does it endure---often for a very long time, capital for centuries, the parasitic bond for decades? A theory of sustainability that could not answer this would be mere moral consolation, predicting a collapse that stubbornly fails to come. The answer is that \emb{the vicious cycle is an open system masquerading as a closed loop.} It sustains itself not by internally regenerating its conditions but by \emph{dumping its entropy on an outside}---a colony, a nature, a future generation, a silenced party---and it defers its own collapse by an ever-widening extraction. Its ``closure'' is space bought with time: so long as there remains a fresh outside to draw upon, the breach in the loop stays hidden, and the system can present itself as self-sustaining while in truth it is being fed from a source it does not acknowledge. This is why capital must perpetually expand---to new markets, to the financialization of nature, to the colonization of the future; expansion is not capital's ambition but its necessity, the continual search for a new outside to consume now that the last is exhausted. And it is why the parasitic bond so often goes hand in hand with the isolation of its victim, the quiet severing of the partner's external supports---so that there is nowhere for the extracted to flee, and no witness to the extraction. \emb{The vicious cycle's endurance is the postponement, and the displacement onto another, of its own collapse.}

Here, then, is the question the antithesis hands to the synthesis. If the cycle is, as the thesis and Eglash alike suggest, all but universal---if generativity must lean on recursion to perpetuate itself at all---and if self-continuation cannot tell good from ill, and if even the return of value cannot exclude the fair stasis or the counterfeit closure, then \emb{wherein lies the difference between the good circle and the vicious one}? The thesis gave us the universality of the cycle; the antithesis has given us its moral indifference, and named the two failures---the stagnant and the counterfeit---that any adequate criterion must exclude. The synthesis must now give us the criterion that the cycle, by itself, withholds: a difference not in the bare fact of self-continuation, which good and vicious share, but in the manner of the return.
   % §3 Antithesis: self-continuation is morally indifferent
% ===== §3 Synthesis: the axiom of sustainability — geometric phase =====
\section{Synthesis: The Axiom of Sustainability and the Criterion of Geometric Phase}
\label{p09:sec:synthesis}

\noindent
The synthesis must supply what the cycle, by itself, withholds: a criterion that distinguishes the good circle from the vicious one without appeal to the bare fact of self-continuation, which both share. I shall argue that the criterion is \emph{geometric}---that the difference between the good and the vicious cycle is, with surprising exactness, the difference between a \emb{circle} and a \emb{spiral}, between a return that accrues nothing and a return that comes home bearing an irreducible remainder. And I shall argue that this is not a metaphor but a structural isomorphism, for which the physics and geometry of \emph{anholonomy}---of the geometric, or Berry--Pancharatnam, phase---supplies the precise form.

\subsection{Return to the root, and the difference between the circle and the spiral}
\label{p09:sec:returnroot}

Begin with the Daoist intuition the prelude already sounded. The \zh{反} of \zh{反者道之动} is not mere oscillation, a swinging back and forth between poles; it is \zh{返}---a \emph{returning}, a coming-home to the root. ``The myriad things in their teeming each return to their root'' (\zh{夫物芸芸，各复归其根}, \emph{Daodejing} ch.~16; \cite{laozi1963}). This returning must be sharply distinguished from a second figure with which it is easily confused: the monotone linear accumulation that Hegel named the \emph{bad infinite} (\emph{schlechte Unendlichkeit})---the mere ``and then, and then,'' the endless additive extrapolation that never comes home but only extends \parencite{hegel2010}. Capital's M--C--M$'$ is bad-infinite in just this sense: it does not return to its root; it extends without rest toward a more that is never enough.

So we already have two figures that must not be conflated: the \emph{returning} of the Dao (return-to-source) and the \emph{accumulating} of the bad infinite (linear extension). But this is not yet the distinction we need, for the antithesis showed that a returning cycle, too, can be vicious (the fair stasis that returns value yet merely repeats). We need a third figure, and it is here that geometry supplies what the two intuitions cannot.

\subsection{The geometric phase: return without return-to-the-same}
\label{p09:sec:geom-phase}

The decisive phenomenon is this. Consider a system carried around a \emph{closed loop} in some space of parameters or situations---carried, that is, so that at the end it returns to the very point from which it set out. One might expect that, having come back to the same point, the system is in the same state. For an important class of systems this expectation is false. When the dynamical (the trivially accumulated) phase is subtracted, the system does \emph{not} return to its initial state: it has accrued a residual---a \emb{geometric phase}---determined purely by the geometry of the path, by the ``area'' the loop encloses, which is to say by the integral of a curvature over the surface the loop bounds. This residual is the phenomenon of \emph{anholonomy}: \emb{the loop closes in the base space, and yet does not close in the fibre.}\footnote{The canonical physical instance is Berry's: a quantum system carried adiabatically around a closed loop in parameter space returns with a phase factor that is not dynamical but geometric---an observable, gauge-invariant quantity equal to the flux of a curvature through the enclosed surface \parencite{berry1984}. The classical optical antecedent is Pancharatnam's \parencite{pancharatnam1956}, for the phase of polarized light carried around a circuit on the Poincaré sphere. What matters for us is the structure, not the particular physics: a closed circuit in the base, an open one in the fibre, the gap between them fixed by an enclosed curvature.}

This is, with a precision that startled the author when it first appeared, the exact form of ``return to the root, yet not to the origin.'' The dialectical movement---thesis, antithesis, synthesis---goes once around and comes back to its word; ``the Buddha speaks of a world, which is then no world, and is thus named a world'' (\zh{如来说世界，则非世界，是名为世界}, the \emph{Diamond Sutra}): the third term returns to ``world,'' to the very word of the first, and yet, having passed through the negation of the negation, it is no longer the naive ``world'' of the outset---it carries a remainder, a sublimation, a geometric phase. \emb{The return happens in the base space (one comes home, to the same person, the same daily round); the sublimation happens in the fibre (one brings home what was not there at setting out).}

We can now state the criterion the synthesis was charged to supply.

\begin{claim}[The geometric criterion of the good and vicious cycle]
Let the state of a relation, or of a value-process, evolve on a fibre bundle: the \emph{base} the space of relational content or situation, the \emph{fibre} the phase of value, of recognition, of jouissance. A cycle of reproduction is \emph{vicious} when its holonomy is zero or negative---when, going once around, it either returns to exactly the point of origin (the stalled oscillation; Hegel's bad infinite is this zero-creating monotone) or suffers a net loss of phase (a net extraction each turn, the rose exhausted). A cycle is \emph{good} when its holonomy is positive---when, in the base, it returns to its root, while in the fibre it has accrued a positive geometric phase. \emb{The circle is the vicious cycle (zero-phase repetition); the spiral is the good cycle (positive holonomy).}
\end{claim}

Note what this criterion accomplishes that the bare notion of self-continuation could not. It \emb{frees ``self-as-end'' from the work of distinguishing good and ill}. Both capital's M--C--M$'$ and the good relational cycle ``take their own continuation as their end''; the difference is not in the autotelic form but in the \emph{mechanism of the phase}. Capital's increment is counterfeited by extraction from an outside (an open system masquerading as a closed loop, in the language of \xref{p09:sec:antithesis}); the good cycle's positive phase is \emph{endogenous}---produced by the curvature of the path itself, by the internal tension-structure of the relation, with no outside drawn upon. This is the geometric form of the distinction between the \emph{stolen} and the \emph{internal} phase that will govern the rest of the paper.

\subsection{The axiom of sustainability: the advance upon generative justice}
\label{p09:sec:axiom}

With the criterion in hand, the paper's central constructive claim---its advance upon Eglash---can be stated.

\begin{claim}[The axiom of sustainability]
Generative justice requires not only that value return to its creators (Eglash's spatial criterion) but that this return suffice, each turn, to reproduce the conditions of generation themselves---and, more strongly, that the cycle accrue a positive geometric phase, a sublimation, with each turn. A cycle that returns value to its creators and yet exhausts the very basis of its own reproduction is just in its distribution and unsustainable in its time. \emb{The axiom of sustainability is the temporal dimension that the spatial criterion lacks: the good circle must return to its root \emph{and} sublimate---it must be a spiral.} This is the paper's contribution to generative justice: not an application of it, but the addition to it of a criterion of time.
\end{claim}

The axiom illuminates, too, why \emph{wu wei}---non-action---will turn out (in \xref{p09:sec:praxis}) to be the fitting practical disposition. If the geometric structure of a system, the curvature of its bundle, is already such that its holonomy is positive, then \emb{the sublimation happens of itself}, and requires no external forcing. \emph{Wu wei} is not inaction but the refusal to extract, to force-close, to impose an external linear end upon a cycle that would, left to its own curvature, sublimate of itself. Forced action---the contrivance of capital, of instrumental reason---is precisely the disease that, in straining to control the phase and wring out an increment, destroys the very curvature by which the system would naturally have produced its sublimation. \emb{Let the cycle be the spiral it would of itself become, and it will sublimate.} This is the geometric--Daoist statement of ``let the circulation of value become a good, natural circle, and it will perpetuate itself of itself.''

\subsection{Two honest caveats}
\label{p09:sec:caveats}

Two debts must be acknowledged at once, for the geometric language is powerful enough to seduce, and the paper had rather name its limits than be caught out at them.

\begin{claim}[First caveat: the phase requires a real bearer]
A geometric phase is physics, and not metaphor, only because the quantity carried around the loop is real---a phase with the structure of a circle (an $S^1$, or a more general group), additive, gauge-invariant, in principle observable. If we speak of ``the phase of value,'' we are obliged to answer: the phase of \emph{what}? What is it that has a periodic, additive, gauge-invariant structure, such that a holonomy can be spoken of at all? Absent a real bearer, the geometric phase is borrowed eloquence and no more. The paper does not discharge this debt here; it incurs it openly, and pays it only in \xref{p09:sec:core}, where the bearer is fixed as the sublimation of \emph{real value}---the jouissance generated and refluxing in the interpretive cycle. Until then the geometric language is to be read as a structural analogy with a forward promise, not a completed formalism. The technical formalization, with its honest treatment of the bearer and the dimensionality, is deferred to the appendix and to the separate Value-Foam study it points to.
\end{claim}

\begin{claim}[Second caveat: ``sublimation $=$ positive phase'' must be argued, not stipulated]
It is not permitted to \emph{define} sublimation as positive holonomy and then announce an isomorphism; that would be to prove by naming. One must characterize, independently, a quantity of ``sublimation'' or ``surplus''---for instance, a margin of sustainability: the net reproduction, each cycle, of the conditions of generation, in excess of what is consumed---and then \emph{show} that, under a suitable bundle formulation, it corresponds to the sign of the holonomy. To weld the two by definition is circular (a vicious circle, in the most literal sense). This paper states the requirement and meets it only in part, in \xref{p09:sec:core}; it flags, here, that the correspondence is a claim to be earned and not a convenience to be assumed. The reader will find, in \xref{p09:sec:turn}, the paper calling itself to account on exactly this point.
\end{claim}

These caveats are not weaknesses to be hidden but the very form of intellectual honesty the paper's epistemology (\xref{p09:sec:epistemology}) will make thematic: that the geometric framework offers a structural \emph{characterization} of what the good circle \emph{is}, and not an operational \emph{procedure} for measuring it from without. To mistake the ``what it is'' for the ``how to measure'' is the positivist's temptation, and the next section is, in part, a guard against it.
    % §4 Synthesis: the axiom of sustainability, geometric phase
% ===== §5 Epistemology: the diagnosis of curvature =====
\section{Epistemology: The Diagnosis of Curvature and the Reversal of the Burden of Proof}
\label{sec:epistemology}

\noindent
The synthesis gave a criterion: the good circle is the one of positive holonomy, the spiral that returns to its root while accruing a remainder. But a criterion of what a thing \emph{is} is not yet a procedure for telling, from outside, \emph{that} it is---and the gap between the two is the whole matter of this section. How, and how far, can the sign of the curvature be read? The section's answer constrains everything that follows, for it arrives at a conclusion that is, at first, disappointing and, on reflection, the very hinge of the paper's practical doctrine: that there is no shortcut to the diagnosis, no external metric that would relieve us of a sustained and fallible attention---and that this absence is not a defect of the theory but a demand of it.

\subsection{The trending-toward-the-good of natural dynamics}
\label{sec:trending}

Begin with an observation that, rightly understood, reverses a burden of proof we scarcely knew we carried. The Daoist tradition offers two homely images, and their homeliness is part of their force, for they record something everyone has seen. A wound left alone will, very often, heal of itself; it is the picking at it, the over-cleaning, the repeated anxious probing, that prevents a self-organizing process wiser than our intervention from doing its work. A body fallen into water will, if it ceases to thrash, float of itself; it is the struggling---the forced action---that defeats the buoyancy by which one would otherwise have risen. \zh{沉疴能自疗，沉溺者不挣则浮}: the deep affliction can heal itself; the one who is drowning, if he does not struggle, floats.

What these images teach is not a vague quietism but something precise about the curvature of certain systems. The wound, the body in water, the convalescing relation: each has a \emph{ground state} that is already a positive curvature, a default tendency toward recovery and self-continuation that operates unless it is interfered with. The body does not need to be taught to float; it floats by its nature, and needs only to stop doing the thing that prevents it. The wound does not need to be made to heal; it heals by its nature, and needs only to be left unpicked. \emb{The default tendency of a natural dynamics is often already toward the good---toward positive holonomy---and such a system goes vicious not for want of intervention but because it has been deformed by intervention.}

This gives \emph{wu wei} a far stronger ground than mere prudence, and the strength of the ground is worth making explicit, for it is easily underestimated. The counsel to ``let be'' is usually heard as a counsel of caution---do less, lest you make things worse---a kind of epistemic humility about our limited foresight. But the claim here is not merely that we are too ignorant to intervene well; it is that the system's own dynamics are, in the relevant cases, already oriented toward the good, so that intervention is not a risky improvement upon a neutral baseline but, characteristically, a deformation of a baseline that was already positive. This reverses a burden of proof:

\begin{claim}[The reversal of the burden of proof]
If the default tendency of natural dynamics is toward the good---if self-continuation and recovery toward positive holonomy are the ground state---then the burden of proof is reversed. It is not \emph{letting be} that requires justification, but \emph{intervening}. The default presumption is that the system will, of itself, tend toward the good; what stands in need of defence is every action upon it. Each instance of forced action must answer a question it is not accustomed to being asked: am I repairing a system genuinely deformed into negative curvature, or am I picking at a wound that would have healed? The modern conviction that a relation ``must be worked at,'' that maintenance is the default and letting-go the deviation, has the burden exactly backwards.
\end{claim}

This is, in fact, the deepest charge the paper lays against the logic of capital and of instrumental reason, and it is worth stating as such. That logic cannot believe that any system will come right of itself; it is constitutionally unable to grant the ground state its positive curvature. And so, by ceaseless management, optimization, and control, it bends every cycle that might have healed into one that now genuinely requires its perpetual management---manufacturing the very dependency it then presents itself as remedying. The relation that ``must be worked at'' becomes, under enough working-at, a relation that can no longer hold itself without the work; the managed system forgets how to float. The reversal of the burden of proof is thus not only an epistemic correction but a diagnosis of how the anxious, managing disposition produces the helplessness it claims to address---the same structure, at the scale of a whole rationality, as the anxious hand that rots the rose's root.

\subsection{Why there is no metric: the honest limit}
\label{sec:no-metric}

But the reversal must not become a licence, and here the section turns sharply upon itself, to forestall an over-strong claim that an earlier formulation of this very argument committed and that must now be retracted. It is tempting to convert the foregoing into a diagnostic rule: that the \emph{amount of forced action a cycle requires to sustain itself} is a reliable marker of its curvature---the good cycle (incentive-compatible, positive-curvature) tending toward self-maintenance with little forcing, the vicious cycle (incentive-incompatible, negative-curvature) demanding ceaseless forcing to prop up its counterfeit closure. There is something to this, and the commercial instance of \xref{sec:praxis} will draw on it; but as a general rule it overreaches, and the overreach must be retracted before it does damage.

The trouble is that one and the same external phenomenon---a high cost of maintenance---answers to at least three quite different internal states, and the external observation cannot, of itself, distinguish them.

\begin{claim}[The maintenance cost is a fallible marker, never a criterion]
The forced action a cycle requires to sustain itself is, at most, an external \emph{marker} of the sign of its curvature, and never a \emph{criterion} of it. One and the same high maintenance cost may answer to three quite different internal states. It may be the building of positive curvature in a founding period---a newly planted rose does need tending, and a young bond does need work, without this impugning its sustainability. It may be the repair of a positive-curvature system under external shock---a sound bond struck by illness or crisis needs upholding, and the upholding is no mark of viciousness. Or it may be the masking of an endogenous negative curvature---the vicious cycle forcing ceaselessly to hide its breach. The external phenomenon \emph{underdetermines} the internal property; the causal inference from marker to property is not deductively valid, and is, in the human case, scarcely susceptible of proof. The maintenance cost is therefore an invitation to diagnose, never the diagnosis itself.
\end{claim}

This limit, properly understood, is not a flaw in the framework but a corollary of its deepest commitment---and the section now makes the corollary thematic, for it governs the whole of the practical doctrine to come. The geometric framework offers a structural \emph{characterization} of what the good circle \emph{is} (a positive holonomy, an endogenous spiral); it does not, and cannot, offer an operational \emph{protocol} for measuring that property from outside.

\begin{claim}[The unobservability of curvature is a demand of the praxis, not a defect of the theory]
If the good and the vicious circle could be told apart by a single external indicator, then ``diagnostic attention''---the sustained, fallible, first-personal labour of reading the curvature---would be superfluous; one could simply consult the metric. But it is precisely the irreplaceability of this attention that will constitute the core of the practice of sustainability. A criterion outsourceable to ``check the maintenance cost'' would itself violate the paper's central claim, for it would settle, by an external sign, what can only be read along a lived path---it would be, at the level of method, a settling failure. The unobservability of the curvature is therefore not an embarrassment the theory must apologize for but a consequence the theory requires: there is no shortcut of judgement, and the sign of the curvature can be read only in a sustained, fallible, first-personal attention, exercised again and again. To mistake the structural ``what it is'' for an operational ``how to measure'' is the positivist's temptation, and the framework's refusal of a metric is its refusal of that temptation.
\end{claim}

So the epistemology arrives, by its own route, at a result that binds the practice to follow, and it is worth drawing the binding explicitly, for it dissolves the usual order of theory and practice. Knowing the curvature is not a measurement but an attention; and because it is an attention---sustained, fallible, first-personal, unoutsourceable---it is not a preliminary to the practice but already the practice itself. \emb{The epistemology is not the prelude to the practice; it is the content of the practice.} To sustain a relation is, in the end, to keep reading its curvature---which is to say, to attend; and the later sections on perception (\xref{sec:perception}) and on the disposition of \emph{wu wei} (\xref{sec:praxis}) will be, in large part, the unfolding of what this attending concretely involves. The attention that the epistemology here identifies as the only access to the curvature is the same attention that the semiotics will recognize as the witness the poem demands, and that the series' earlier papers described as joint attention. There is, in this convergence, a first intimation of the paper's recurring discovery: that to know a relation rightly and to sustain it rightly are not two acts but one, and that both are forms of an attention that declines to possess what it attends to.
 % §5 Epistemology: the diagnosis of curvature
% ===== §5 Praxis: the strategy of no-strategy =====
\section{Praxis: The Strategy of No-Strategy}
\label{p09:sec:praxis}

\noindent
The epistemology ended on a result that dictates the form of the practice: that sustaining a relation is a sustained, fallible, first-personal attention, not the application of a metric. This section draws out what such a practice can and cannot be. Its conclusion is, at first, disquieting---that there is, strictly, no \emph{strategy} for sustainability---but the disquiet is the doorway to the answer, and the answer is the one the whole paper has been circling: a \emb{strategy of no-strategy}, which is \emph{wu wei} given its precise and non-mystical content.

\subsection{Why sublimation cannot be aimed at}
\label{p09:sec:no-aim}

Lay the foregoing in a line and an uncomfortable necessity appears. The mark of the good circle is positive holonomy---sublimation; and holonomy is accrued along a path walked in the real, not cashed out by a sign (\xref{p09:sec:semiotics} will make this the heart of the semiotics). But any intervention that takes ``producing sublimation'' as its direct aim is attempting to produce, directly, the very thing that can only arise as the by-product of a walked path---which is the definition of settling failure. \emb{Any strategy that takes the production of sublimation as its direct end will, precisely thereby, destroy it.} The more one schemes to ``manage the relation into sustainability,'' the more one flattens the spiral with the instrumental reason---the control, the extraction, the direct wringing-out of phase---that the geometric criterion identifies as vicious. A strategic love is not love; it is settling. To aim at the holonomy is, in the geometry, a category error: one cannot ``walk toward'' a holonomy, one can only walk the path, and the phase accrues of itself.

\begin{claim}[The strategy of no-strategy]
Sustainability has no strategy, only a disposition. The good circle cannot be produced, only not-obstructed. Every ``technique for making a relation sustainable'' is, structurally, a forced action, an extraction, and destroys what it would preserve by cashing out a sublimation never walked. The only self-consistent practice is a negative, custodial one: not to do something so as to produce sublimation, but to refrain from doing what destroys the curvature under which sublimation arises of itself---not to extract, not to hasten, not to settle, not to counterfeit real presence with a cheap sign. \emb{\emph{Wu wei} is not the absence of practice but the redirection of practice from ``producing sublimation'' to ``guarding the curvature under which sublimation occurs.''}
\end{claim}

\subsection{Three domains, three insights}
\label{p09:sec:three-domains}

The disposition shows itself, with a different lesson, in three domains; their parallel is part of the evidence that the structure is genuinely universal and not forced from a single case.

\emc{The Daoist domain} contributes the ground: that the natural dynamics trends toward the good (\xref{p09:sec:trending}), so that letting-be is not a passive default but a trust grounded in the system's own positive curvature. The wound heals; the body floats; the relation recovers.

\emc{The commercial domain} contributes a sharpening, and a caution. The highest selling is, very often, not to sell. To press a sale is to emit a self-undermining signal: the very act of pressing betrays that ``my interest and yours diverge, and so I must persuade you''---and the harder one presses, the more loudly one announces that the exchange may not, in truth, serve the buyer. The credibility of a relation of exchange cannot, in the end, rest on the polish of the pitch (which is cheap, outsourceable, machine-generable); it rests on the unforgeable fact that the thing is genuinely good, the relation genuinely mutual---so that the customer \emph{comes of themselves}, and a mutually beneficial circulation arises that needs no forcing. Incentive-compatibility produces an endogenous positive curvature: when the parties' interests align in the cycle, the cycle perpetuates itself with little forcing, and \emph{wu wei}---not interfering with a mutually beneficial cycle that would continue of itself---is the fitting disposition. (The maintenance cost reappears here as a marker of curvature; but it is held, as \xref{p09:sec:no-metric} demanded, to the status of a fallible marker, never a criterion.)

\emc{The domain of the rose} contributes the sharpest lesson of all, and lifts the doctrine from ``forced action is superfluous'' to ``forced action is harmful.'' One who, dreading that the rose will wilt, waters it without cease, changes its water, tends it with an anxious love---that one cannot alter the fact that the flower must fall, and worse: he rots the root, and brings about the withering he feared. Here forced action does not merely fail to help; it \emph{manufactures the very outcome it dreads}. The fear of an end, made into action, produces the end. This is \zh{为者败之，执者失之}---``the one who acts upon it ruins it, the one who grasps it loses it'' (\emph{Daodejing} ch.~64)---and it cuts at once through capital (whose dread of the end of growth drives the extraction that exhausts its base), through the intimate bond (whose dread of loss drives the control that rots the relation), and through the Dao itself. \emb{The deepest unsustainability is born of the fear of unsustainability.}

\subsection{Wu wei is not negligence: the diagnostic anchor}
\label{p09:sec:not-negligence}

But \emph{wu wei} must be sharply distinguished from negligence, or the doctrine becomes a high-toned alibi for abandonment---and the Daoist tradition has, historically, been put to just this use, sanctifying a ruler's inaction as cosmic wisdom. The paper must therefore furnish a criterion that divides the two; and the criterion is, exactly, the sign of the curvature.

\begin{claim}[Wu wei versus negligence]
\emph{Wu wei} is non-intervention in a positive-curvature system; negligence is non-intervention in a negative-curvature one. The same ``doing nothing'' guards a sublimation arising of itself in the first case, and tolerates an extraction underway in the second. Whether one ought, at this moment, to let be or to act is not settled by the feel of one's own disposition but by the diagnosis of the curvature: is this cycle, now, sublimating (returning to its root while the fibre accrues a positive phase), or is it being exhausted (a net extraction, one party giving without cease and nothing returning)? The legitimacy of \emph{wu wei} is local, and stands in need of ceaseless re-diagnosis; it is not a posture adopted once for all. The root of the rose may be left to its own regeneration; the partner being drained dry may not be left to ``wu wei'' her way to depletion---there, what is called for is a repairing action.
\end{claim}

\subsection{Repair is subtraction}
\label{p09:sec:repair}

What, then, is the repairing action that a negative-curvature system calls for? Not a new forced action laid on top of the old, but the \emph{removal} of the forced action that is doing the harm. This is the reconciliation of repair with \emph{wu wei}, and it is the section's final claim.

\begin{claim}[All true repair is subtraction]
Repair is not the addition of a new intervention but the subtraction of the one that sickens. The drowning body does not need to be taught a more elaborate stroke; it needs to be brought to cease its thrashing. The wound does not need a fiercer medicine; it needs the picking hand withdrawn. The intervention a negative-curvature system requires is, at bottom, the withdrawal of the anxious, extracting, force-closing hand---so that the system's own positive curvature may resume its work. \emb{The highest repair is itself a kind of \emph{wu wei}: not to do more, but to take away the hand that has all along been picking at the wound.} So \zh{扶}---to ``steady'' the drowning---does not mean to seize and haul him up (that is forced action, complicit in the thrashing); it means to hold him so that he ceases to struggle, and the water bears him up of itself. To steady is the least action that makes way for self-healing.
\end{claim}

The two dispositions are thus unified. \emph{Wu wei}, in a positive-curvature system, guards the sublimation arising of itself; repair, in a negative-curvature one, withdraws the hand that has bent the curvature negative---and thereby restores the conditions under which \emph{wu wei} is once more legitimate. Repair completes itself by giving way again to non-action. The practice of sustainability is, in sum, neither management nor abandonment but a vigilant, custodial attention: to keep reading the curvature, to stand back where it sublimates of itself, to withdraw the harming hand where it is being drained, and never to try to cash out, with a sign, the sublimation that can only be walked. It is neither a strategy nor a laissez-faire, but a poetic custody---like the reading of a poem that is never read out.
       % §6 Praxis: the strategy of no-strategy
% ===== §6 Semiotics: what the sign can and cannot do; the poem as paradigm =====
\section{What the Sign Can and Cannot Do: The Poem as Paradigm of the Good Circle}
\label{p09:sec:semiotics}

The preceding sections erected a criterion. The good circle is a spiral and not a circle: it returns to its root while accruing an irreducible remainder, a surplus we have named, in the language of geometric phase, a positive holonomy. The sections on epistemology and praxis then asked what finite agents can do in the face of such a criterion, and answered with a disposition rather than a technique---not to \emph{produce} sublimation but to refrain from destroying the curvature under which sublimation arises of itself. Yet that answer left a promissory note unredeemed. For whether we speak of guarding or of repairing, we must \emph{do} something; and everything we can do must pass through a single interface---the \emb{sign}.

This section takes up that interface. It must answer three questions, and the order is dictated by a principle that has governed the series throughout: epistemology before praxis. \emph{Why} can a sign alter the trajectory of a cycle at all (\xref{p09:sec:two-mechanisms})? \emph{When and how} must a sign intervene if it is to bend the cycle toward the good (a question that begins here and is completed only in \xref{p09:sec:core})? And \emph{when and how} does a sign fail (\xref{p09:sec:failures})? The section's summit (\xref{p09:sec:poem}) will then advance what may be among the strongest claims of the paper: that among all signs, the \emb{poem} occupies a singular place---that it is the pure form of the good circle's sign, the Archimedean point of the entire semiotics.

\subsection{The paradox of the sign: the sole operable thing, and the gateway to alienation}
\label{p09:sec:symbol-paradox}

We must begin from a tension, and resist the temptation to dissolve it prematurely; the tension is itself the engine of the section.

On the one hand, \emb{the sign is the only thing we can formally operate, know, and transmit.} The real of love---that connection between two beings which is unsayable, immeasurable, exhausted by no concept---is not, in itself, operable. We cannot ``adjust'' it directly, cannot hold it in the hand and add to or subtract from it. What we can lay hands on is only the sign: a declaration, a gift, a rite, a line set down on paper. All practice must pass through this narrow gate. In this sense the sign is not the ornament of a relation but the sole stratum at which a relation can be deliberately entered. To deny the sign its office is to deny the very possibility of practice.

On the other hand---and this is no less true---\emb{the real of love does not reside in signs.} If a relation subsists on signs alone, it has already stepped into alienation: the signifier slides, multiplies, and is exchanged, yet anchors no real. This has been a constant concern of the series, and it acquires today an unprecedented edge. Technology, and above all the advent of large language models, is driving \emb{the marginal cost of the signifier toward zero}. A machine can generate an inexhaustible supply of endearments, of bespoke declarations, of impeccably calibrated phrasings for a gift. When signs can be produced this cheaply and at this scale, the sign's function as a \emph{credible signal of investment and commitment} collapses. This is not a distant worry but an inflation of the sign already underway.

Set the two faces side by side and the genuine problem of a semiotics comes into view; and it hangs, precisely, upon the geometric framework erected above.

\begin{claim}[The paradox of the sign]
The sign is the sole operable interface upon the dynamics of a relation, yet the sign itself \emph{carries no} geometric phase---phase is accrued along a path, and a path is a process walked through in the real, incompressible. The sign can \emph{initiate, redirect, mark} a cycle; it cannot \emph{substitute} for the cycle's own unfolding. \emph{Alienation} is precisely the error of taking the throwing of a sign for the walking of a path---the attempt to cash out, directly, by means of the sign, a sublimation that ought to have arisen as the by-product of a path actually traversed. And artificial generation has industrialized this very illusion: by rendering ``the throwing of a perfect sign'' all but free, it tempts us, as never before, to counterfeit the real path with the sign.
\end{claim}

This is why ``the sign growing ever cheaper'' is a threat at the core of sustainability---not because the sign has ceased to matter but, on the contrary, because the cheap sign lets the vicious cycle (the cycle of zero area, of immediacy, of the skipped process) masquerade as the good one. A love-letter ghost-written by a machine, flawless in its diction, and a clumsy sentence spoken aloud to one's face, freighted with real hesitation---on the ``quality'' of the sign the former far surpasses the latter; but on the real path each anchors, the latter may carry the whole of the phase while the former tends toward zero. The entire unfolding of this section is a reverberation of this paradox.

\subsection{How a sign alters dynamics: two mechanisms}
\label{p09:sec:two-mechanisms}

To answer ``when should a sign intervene'' one must first answer ``by what power can it intervene at all.'' This is the order of epistemology before praxis: only by grasping the mechanism does practice escape blind trial. I claim that the sign alters the trajectory of a cycle by two \emph{independent} mechanisms, belonging respectively to psychoanalysis and to political economy. They alter the dynamics differently, and must be set out separately before their synthesis can be seen.

\subsubsection{The psychoanalytic mechanism: the sign anchors the topology of the loop}

The first mechanism is structural. In the Lacanian register the sliding of the signifier is in principle endless---meaning defers without rest along the chain, with no natural mooring. A signifying act---an ``I love you,'' a naming, a rite---functions as a \emb{quilting point} (\emph{point de capiton}): retroactively it pins down the meaning that until then floated free, producing a stable signified.\footnote{Lacan's figure is the upholsterer's button that fastens the stuffing: it is the discrete anchors that lend structure to an otherwise loosely sliding fabric. The crux is that the anchoring is \emph{retroactive}---the meaning does not wait, in advance of the sign, to be designated; it is constituted retroactively at the moment the sign falls \parencite{lacan2006}.} 

Restated in the language of dynamics: the sign does not inject energy into the system; it \emph{alters the system's topology}. It manufactures, in the phase space of the relation, a new fixed point or attractor, or---and this is the decisive one for us---it \emb{closes an otherwise open, dissipative trajectory into a loop}. A declaration can bend the course of a relation not, as a rule, because it conveys new information (both parties most often already know), but because it anchors the real diffused between them, as yet unformed, into a structure that can be \emph{repeated, and therefore reproduced}.

This matters for the whole paper, because \emb{without a loop there is no holonomy to speak of}. The geometric phase is a path-integral around a closed loop; and for there to be a closed loop, something must first suture the open trajectory. The sign is that suture. Hence the first epistemic office of the sign:

\begin{claim}[The topological office of the sign]
The sign structures the unformalizable real into a fibre bundle capable of bearing a geometric phase. The sign does not produce sublimation, but it delineates the \emph{loop} within which sublimation can occur. This is the first ground of what the sign \emph{can} do, and the reason it cannot be dispensed with: absent the anchoring of the sign, the real of a relation is merely a dissipative flow, with no remainder to accrue.
\end{claim}

\subsubsection{The political-economic mechanism: the sign coordinates the good equilibrium}

The second mechanism is selective, and a different vocabulary is needed for it. The sustainability of a relation---of any cooperative system---is often a problem of \emb{multiple equilibria}. There exists in the system a ``good'' equilibrium in which the two parties mutually invest and value circulates back to each, and a ``bad'' equilibrium in which they mutually defend and each extracts. Which the system falls into is determined not by who is more virtuous but by \emph{the parties' shared belief about one another's intentions}---common knowledge, in the game-theoretic sense.

The sign---and above all the costly, public, irreversible sign---here functions as a \emb{coordinating device}. A weighty gift, a public rite, a covenant written down and witnessed: their office is to manufacture common knowledge---``we both know, and each knows the other knows, that the game we are playing is the good one.'' It is this layer of shared belief that pushes the system out of the basin of the bad equilibrium and into the basin of the good. This explains a fact long obscured by romantic talk: why the gift must be \emph{costly}, the declaration \emph{public}, the covenant \emph{written and witnessed}---cost and publicity are not ornament but the credibility of the signal, and credibility produces shared belief, and shared belief flips the equilibrium.\footnote{This is the logic of costly signalling \parencite{spence1973,zahavi1975} and of coordination by focal points \parencite{schelling1960}. But placed within the framework of generative justice it acquires a new sense: the political-economic office of the sign lies not in ``conveying information'' but in serving as the focal point that makes the good equilibrium common knowledge.}

\subsubsection{The synthesis of the two mechanisms}

The two mechanisms can now be joined, and the synthesis is itself a new claim. The psychoanalytic mechanism says: the sign \emph{constructs the topology of the loop}, making sublimation structurally possible. The political-economic mechanism says: the sign \emph{selects the basin of the loop}, making the good circle the one jointly chosen. One is topology; the other is selection.

\begin{claim}[The double office of the sign upon dynamics]
The sign alters the dynamics of a relation at two levels: it \emph{delineates} the structure of the phase space (psychoanalysis: anchoring the loop), and it \emph{coordinates}, among multiple stable states, the good branch (political economy: selecting the basin). The sign-inflation of the age of artificial generation attacks both at once---it devalues the anchor (signifiers proliferate and can no longer pin a real), and it strips the signal of its cost (no longer credible, it can coordinate no common knowledge). The crisis of the sign is therefore a \emph{double} crisis.
\end{claim}

Note that the two mechanisms do not reduce to one another. The topological anchoring can take place within a single mind in solitude (a vow one speaks to oneself anchors one's own real), whereas the coordination of an equilibrium is essentially multi-agent. They are two incommensurable faces of the sign's power, and they prefigure the polyphony the paper's final section will avow: one and the same phenomenon shows itself, at different epistemic stations, as different and mutually irreducible truths.

\subsection{How a sign fails: three pathologies}
\label{p09:sec:failures}

Having grasped what the sign \emph{can} do, we can characterize precisely what it cannot. The failure of the sign is not single but takes three structurally distinct modes. To distinguish them is essential, for each answers to a different pathology and demands a different response.

\subsubsection{Inflationary failure: the surfeit of the sign}

The first failure is the \emb{surfeit} of the sign. The signifier multiplies without bound while the real investment of the signified does not grow with it. The contemporary performance of relation---the staged tenderness of social media, the algorithmically suggested phrasing of an anniversary, the machine-generable endearment---is its type. Here both offices of the sign fail at once: anchoring fails (too many signifiers, none can pin a real), and signalling loses faith (cost at zero, no common knowledge can be coordinated). The cycle seems busier, more ``full of content,'' than ever, and yet it idles---its geometric phase tends to zero. This is the dominant pathology of the present, and the direct issue of the sign's inflation.

\subsubsection{Settling failure: the usurpation by the sign}

The second failure is more hidden and more profound---and it is precisely the pivot of this section toward the poem. Here the sign is used to \emph{substitute} for, rather than \emph{invite}, a path: to settle with a sentence a wound that needed time to heal, to discharge with a gift a debt that needed presence to be borne, to ``complete'' with a perfect phrase a process that ought to have been walked. The sign here is no longer the quilting button that anchors the loop; it \emb{declares the loop already closed}---though in fact it is not.

\begin{claim}[Settling failure]
Settling failure is the use of a sign to \emph{cash out} a sublimation never walked through. It counterfeits, as a balance the sign can draw on at once, a phase that ought to have arisen as the by-product of a path. Its essence is fraud---against the other, and against oneself. For phase depends on the path and not on the endpoints: two relations may arrive at ``the same situation,'' yet, having walked different paths, accrue different sublimations. Sublimation cannot be acquired instantaneously; it can only be walked. Every sign that proclaims ``we have arrived'' while in truth it never set out is a settling failure.
\end{claim}

It is this failure that forces an unintuitive insight: \emb{depth admits no shortcut}. In the language of geometric phase, the accrued phase equals the area enclosed by the path---the integral of the curvature around the loop. Hence the phase depends on the path and not on the endpoints. Carried over to relation: what we call ``depth'' just \emph{is} the area the path encloses. One cannot leap over the process and possess a deep relation directly---depth \emph{is} that incompressible area. So ``instant intimacy,'' ``express devotion,'' must be, structurally, vicious cycles: they shrink the area of the loop to the minimum (shortest path, zero waiting), and the geometric phase tends to zero, and no sublimation occurs. The ``immediacy,'' ``acceleration,'' ``efficiency'' extolled by capital's logic are precisely the systematic annihilation of geometric phase---the flattening of a spiral that might have climbed into an oscillation in place. \emb{Slowness, therefore, is a formal necessary condition of the good circle, and not an aesthetic preference.}

\subsubsection{Ossifying failure: the petrification of the sign}

The third failure is \emb{petrification}. A sign that once anchored a good circle hardens into an empty rite that no longer connects to any real---the perfunctory ``I love you,'' the commemoration become mere husk, the gift become duty. The loop is still there, but it is now the pure repetition of zero curvature (what the earlier sections called the ``zombie cycle that refuses to wilt''). The sign here is not at fault; the fault is that it is no longer re-walked in the real; it has decayed from anchor into inertia.

The remedy for ossifying failure echoes, exactly, the ``let it wilt and then regenerate'' of the sense of \emph{wu wei}: to let the old sign die, and let a new anchoring grow back out of the real. To try to rescue a petrified sign by more frequent, more forceful repetition only hastens its rigor mortis---which is the structure recurring throughout (and in the earlier praxis section): the action born of fear manufactures the very outcome it fears.

\subsection{The poem as paradigm of the good circle's sign: the Archimedean point of the semiotics}
\label{p09:sec:poem}

With the three failures before us, a naive conclusion seems to suggest itself: since the sign is so dangerous, let us use fewer signs, return to some presemiotic innocence. This conclusion is false, and instructively so. For \emb{the contrary of settling failure is not ``fewer signs'' but a special kind of sign---one that refuses to settle itself}. And among all signs, the one that brings this ``refusal to settle'' to purity is the \emb{poem}.

This is the core claim I wish to make: the poem is not one example within the semiotics but its \emph{Archimedean point}---a fulcrum from which the whole understanding of ``what a good sign is'' may be levered. Let me unfold, in three layers, the poem's title to being the paradigm of the good circle's sign.

\subsubsection{First layer: the poem structurally refuses closure, and is therefore immune to settling}

A contract, a sentence reading ``the matter is resolved,'' a machine-generated standard endearment---their semiotic function is \emph{termination}: to declare the loop closed, the meaning delivered, the file ready for archiving. The poem cannot do this, and \emph{refuses} to. Every time a poem is read, its meaning unfolds anew; it is never ``read out.''\footnote{Here ``poem'' does not designate the literary genre of lineated text but a manner in which a sign operates---which may equally take place in a passage of prose, a rite, a held gaze, and may fail to take place at all in a formal ``poem'' (see the false poem of \xref{p09:sec:false-poem}).} Precisely because the poem cannot be settled, it cannot be used as the instrument of ``cashing out an unwalked sublimation.'' In other words, \emb{the poem is, in its very form, immune to settling failure.} This is its first title to being a paradigm: it is a self-defending sign, a sign that structurally refuses to be put to usurping use.

\subsubsection{Second layer: the poem moves the site of sublimation from the text to the path of interpretation}

This is the most important layer, and the interface that joins the poem to the geometric framework. The meaning of a poem does not lie in the arrangement of its signifiers---this is why we say ``the formal content of the poem matters less''; the meaning of a poem lies in \emb{its unfolding within the dynamics of the interpreter}.

In the language of geometric phase: the text of the poem is only a loop (or a point of departure) in the base space of situations; meaning, sublimation, is \emph{the holonomy the interpreter accrues by traversing that loop once around}.

\begin{claim}[The meaning of a poem is the holonomy of the interpretive path]
The meaning of a poem $=$ the area the interpreter encloses along its loop. The same poem, for different interpreters---and for the same interpreter at different moments---walks different paths, encloses different areas, accrues different phases. This is why a poem is ``inexhaustible'': its loop permits an area arbitrarily large, and unrepeating, to be enclosed. Each rereading is a spiral and not a circle---it carries home a remainder the last did not.
\end{claim}

From this follows a surprisingly operable measure of ``poeticity'': \emb{the poeticity of a sign equals the supremum of the non-trivial holonomy it can bear.} A slogan, an advertising jingle---its loop is too small, too closed; on a second recitation the phase is already zero; it has been ``read out''; its poeticity tends to zero. A great poem withstands a lifetime of rereading precisely because its loop permits the accrual of an almost unbounded, unrepeating phase. Poeticity is not a vague term of praise but a structural magnitude---a magnitude concerning the capacity of a sign to bear sublimation.

\subsubsection{Third layer: the poem exhibits the epistemic structure of witness}

The third layer fills the gap left by the first mechanism (anchoring). The poem does not transmit to the interpreter a ready-made signified; it \emb{invites the interpreter to walk a real once over in person, and to become, in the walking, a witness.} This is the pure form of that act by which a sign ``anchors the loop without substituting for the path''---the poem, with the fewest signifiers, sets in motion the largest unfolding of the real within the interpreter's own dynamics.

And every sign-intervention done rightly---a declaration, a gift, a covenant---does the very same thing: it does not ``inform'' the other of love; it \emph{invites} the other to walk love once over, in person, within their own dynamics, and there to become a witness. The poem, then, is not one example within the semiotics but \emb{the limit form of all legitimate sign-intervention}---the ``invitation of the path'' brought to purity, the ``settling'' reduced to zero.

\begin{claim}[The poem as the apex of the good sign]
Every good sign-intervention, at the limit, approaches the structure of the poem: it \emph{anchors} the loop while refusing to \emph{close} it, it \emph{initiates} sublimation while refusing to \emph{cash it out}, it \emph{invites} the interpreter to walk the real in person and there to witness generativity itself. Settling failure is precisely the moment a sign betrays its own poeticity and usurps the office of a contract. The crisis of the sign in the age of artificial generation is, in essence, \emph{the poeticity of the sign being devoured by its settleability}---the infinitely generable signifier tempting us to treat every sign as a settleable deliverable, and thereby annihilating holonomy. Hence the practice of resisting alienation just is the \emph{restoration of the sign's poeticity}: to let the sign be, once more, an invitation to a real path, and not the cashing-out of a sublimation.
\end{claim}

This claim carries, besides, a retroactive dividend. The eighth paper of this series treated the lyric's address to the big Other---the poem presents generativity itself, rather than any ready-made structure. There that thesis concerned the lyric; here it is re-sited: the structure of the lyric just \emph{is} the semiotic paradigm of the sustainable good circle. The series, then, is not eight papers and then one more, but a ninth that \emph{retroactively re-anchors} the meaning of the eighth---which is itself a \emph{point de capiton}, a spiral return. \emb{The series enacts, in its form, the content it asserts}: no single paper settles the meaning, which goes on generating new remainders along the path of the papers' mutual rereading.

\subsection{The vicious twin of poeticity: the false poem}
\label{p09:sec:false-poem}

But we must at once set a guard upon this too-beautiful picture, lest it collapse into a romanticization of ``poetic vagueness.'' That the poem refuses to settle is its strength; yet \emb{``never settling'' may slide into another pathology}---an infinite deferral, a game of the signifier that never redeems itself: the pure sliding of signifiers, the cynic's irony, the posture that pleads ``staying open'' as an alibi for never investing. This is in truth a higher form of inflationary failure, wearing the mask of the poem.

Poeticity, then, has its vicious twin. I call it the \emb{false poem}. The difference between the true and the false poem returns, in the end, to geometric phase:

\begin{claim}[The criterion of true and false poem]
The \emph{true} poem's non-closure is in order that sublimation \emph{be lived}---the interpreter walks the path in person and accrues a positive holonomy. The \emph{false} poem's non-closure is in order to \emph{evade} the cost of sublimation---its holonomy is zero, it proclaims ``openness'' in place while no one has truly set out. \emph{Openness} is not, in itself, the good; only the openness that bears a positive phase is. A true poem invites a real path to be walked; a false poem pleads ``openness'' as an alibi for never investing.
\end{claim}

This distinction is crucial, for without it the section would be justly charged with aestheticizing vagueness, with dressing ``nothing is certain'' as profundity. The false poem's refusal of closure is an idling refusal---it never demands of the interpreter the cost of walking a real in person, and so it breeds countless postures and zero sublimation. The true poem's refusal of closure is a generous refusal---it hands the whole weight of sublimation back to the path the interpreter lives. On the surface both ``give no ready answer''; in holonomy they are opposites.

And here the question of the \emph{bearer} of the phase can no longer be deferred: that ``generativity'' the poem presents, that ``existence'' the interpreter witnesses---of what is it the phase? We have twice now arrived at the threshold of one and the same answer---once following the logic of the sign's anchoring, once following the logic of the poem's unfolding---and both independent threads point to a single bearer: \emb{the sublimation of jouissance} \parencite{lacan1998}. The poem is the pure artistic form of this sublimation.\footnote{In the Seminar on \emph{The Ethics of Psychoanalysis} \parencite{lacan1992} Lacan expounds sublimation precisely by way of the relation between the poem and \emph{das Ding}: sublimation is not the satisfaction of desire but the ``raising of an ordinary object to the dignity of the Thing''---encircling that impossible, empty real without filling it. This structure will become, in \xref{p09:sec:core}, the heart of the definition of the appropriate poem.} But to fix this bearer securely, and to say what divides the ``good'' poem from the ``vicious'' one, we must first pass a threshold---the next section shows that poeticity, in itself, \emph{does not guarantee} the good.
    % §7 Semiotics: the poem as paradigm
% ===== §7 Poeticity does not guarantee the good =====
\section{Poeticity Does Not Guarantee the Good}
\label{p09:sec:turn}

The preceding section raised the poem to a great height---the paradigm of the good sign, the Archimedean point of the semiotics. Such praise, if it is not at once and severely qualified, would bring the whole paper down here. For an objection waits at the door, and it is sharp enough to be fatal: \emb{poeticity guarantees only the self-continuation of generativity, not the good.}

The objection is fatal because it turns the paper's own blade against the paper. The \emph{antithesis} of this paper (\xref{p09:sec:antithesis}) laid down an iron law: \emb{self-continuation, in itself, is morally indifferent.} Capital's accumulation continues itself; the rose's blooming continues itself; the parasitic bond continues itself---``ceaseless circulation,'' alone, never constitutes the good. We must now concede: \emph{the poem is no exception.} A generator of $\lambda>1$, bearing a phase, unfurling without end---by what right is what it generates the good circle? By what right does the poetic spiral necessarily climb, rather than descend?

To speak honestly: there is no such guarantee. And the counter-examples lie ready to hand, and must be faced.

\subsection{The vicious poem exists}
\label{p09:sec:vicious-poem}

\begin{claim}[Poeticity does not entail the good]
There exist genuinely poetic, yet vicious, generators of generativity. Fascist aesthetics is poetic---it is by no means a dead grammar; it installs in its followers a grammar of meaning able to go on self-generating, self-``sublimating,'' impassioned, unending, demanding personal witness, refusing to be settled. Cultic speech is poetic. The addict's inner monologue is poetic---it ceaselessly generates new rationalizations, new cravings, self-reproducing, never closing. All satisfy the criterion of poeticity given in the preceding section: $\lambda>1$, the interpreter personally drawn in, refusal of settling, spiral and not mere repetition. They are \emph{genuine} generators of generativity, and yet vicious. Therefore ``poeticity $\Rightarrow$ good'' is a false proposition.
\end{claim}

We must pause and see exactly where the preceding section erred---for the error is the very one the paper had earlier warned itself to guard against. In \xref{p09:sec:synthesis} we set down a candid caveat: the correspondence ``sublimation $=$ positive phase'' requires an argument independent of definition, on pain of circularity. And \xref{p09:sec:semiotics}, in praising the poem, did precisely this---it quietly stitched ``poeticity'' to ``the good'' by the sign of the holonomy, without ever arguing, independently, \emph{why the poetic phase must be positive}. This section comes to discharge that debt.

The manner of discharge is to concede a stratification:

\begin{claim}[Poeticity is necessary but not sufficient for the good]
Poeticity guarantees the \emph{form} of the cycle---that it is a living spiral, not a dead circle (ossification) nor a settled closure (the contract). Poeticity guarantees that the cycle \emph{can} endure, \emph{can} sublimate. But it does \emph{not} determine the \emph{direction} of that sublimation. Poeticity gives the cycle the capacity for endurance, but not the direction of endurance. The good must be guaranteed by something else.
\end{claim}

\subsection{Two orthogonal criteria: poeticity and justice}
\label{p09:sec:orthogonal}

What is that ``something else''? To pursue it is to drive the paper's two languages to the point at which they must separate, and then be rejoined. I claim:

\begin{claim}[The orthogonality of poeticity and justice]
\emph{Poeticity} is a \emph{dynamical / formal} criterion: it asks---can this cycle endure? is it a spiral or has it gone rigid? \emph{Goodness} is a \emph{justice} criterion: it asks---is the phase of this cycle drawn from an internal reflux, or extracted from an outside? whose conditions of generation does it reproduce, and whom does it reduce to expendable fuel? The two criteria are \emph{orthogonal}: a cycle may be poetic and vicious (fascist aesthetics), or just and rigid (a bond fair but long since lifeless). The formal ``aliveness'' and the just ``goodness'' are two independent dimensions.
\end{claim}

With this orthogonal decomposition, fascist aesthetics can be diagnosed precisely, and no longer merely reviled as ``bad.'' It \emph{is} poetic---it meets every formal condition; it is \emph{not} just---its spiral can go on sublimating only because it takes ``the enemy,'' ``the alien,'' ``the other'' as expendable fuel. Its inner ``sublimation'' is \emb{stolen}: the phase is produced not by the tension among the participants internal to the cycle but extracted from an outside excluded and consumed, and then disguised as an internal sublimation.

This drives an earlier insight to its terminus. In \xref{p09:sec:antithesis} we said that the vicious cycle is ``an open system masquerading as a closed loop''---it dumps its entropy on an outside, deferring its own collapse by an ever-widening extraction. We now see that this structure has, at the level of poeticity, its exact twin:

\begin{claim}[The vicious poem is the stolen phase]
\emph{Fascist aesthetics is to the poem as capital is to generativity}---formal twins, ethical opposites. Capital ``self-augments,'' formally resembling generativity, but its increment is an extracted, counterfeit increment; the vicious poem ``self-sublimates,'' formally resembling the true poem, but its phase is a stolen, counterfeit phase. Both counterfeit a remainder that ought to have been internal, by excluding certain beings as pure fuel. A cycle's poeticity answers ``can it survive''; its justice answers ``\emph{should} it survive, and \emph{on whose blood} does it survive.''
\end{claim}

\subsection{An interrogable criterion: who is the pure fuel?}
\label{p09:sec:fuel}

The orthogonal decomposition gives us a concrete, reiterable diagnostic question---one that will become the point of departure for the next section's constructive work. How to tell the internal sublimation from the stolen one? Look whether the cycle has, within it, an outside systematically reduced to \emb{pure fuel}---a being that only contributes and never benefits: a silenced partner, a colonized other, a group consumed as an enemy, a future overdrawn, a nature exhausted.

\begin{claim}[The pure-fuel criterion]
Whether a spiral's sublimation is good or vicious may be interrogated by this question: \emph{is there, within this cycle, some one person or thing that exists only as cost---excluded from benefiting in the very sublimation it feeds?} If there is, this is a vampiric poeticity---its phase is stolen. In the good poetic cycle there is no ``pure outside'': every being drawn in is both producer and beneficiary of the phase; the sublimation refluxes to every one who generates it.
\end{claim}

This criterion is the ethical redemption of the ``open system masquerading as a closed loop'' of \xref{p09:sec:antithesis}, and it carries the core principle of generative justice---that value, unalienated, should return to those who create it---over from the domain of matter and labour into the domain of sublimation and witness. It will be locked, in the next section, into the substantive content of the ``appropriate'' poem.

It must finally be stressed where this section stands in the logic of the whole. To concede that ``poeticity does not guarantee the good'' does not weaken the preceding section; it is the very precondition of that section's standing. For it hands ``the good'' back to an independent, just criterion, and thereby forbids the poem---and the geometric form behind it---to usurp the role of a meta-framework that subsumes all the rest. The poetic framework can adjudicate ``sustainability / endurance''; but to adjudicate ``the good'' it must hand the microphone to the framework of justice. \emb{The poem answers whether the cycle can live; justice answers whether it should, and on whose blood.} These two questions are orthogonal, and therefore require two mutually unsubsuming frameworks---the polyphonic structure the paper's final section will avow, here for the first time appearing in a form impossible to evade: that even the most beautiful of things cannot, alone, vouch for the good.
         % §8 Poeticity does not guarantee the good
% ===== §8 The appropriate poem: the non-possessive interpretive cycle that generates real value =====
\section{The Appropriate Poem: The Non-Possessive Interpretive Cycle That Generates Real Value}
\label{sec:core}

\noindent
Up to here the greater part of the paper has been spent diagnosing the vicious---failure, extraction, the stolen phase, the outside reduced to fuel. This was a necessary detour, but the time has come to put the question affirmatively: since poeticity does not guarantee the good, \emb{what is the appropriate, the just, the self-perpetuating poeticity}?

This section is the keystone of the paper. Its particularity is this: the answer is not derived along a single line, but is the meeting, at one and the same point, of \emph{two independent inquiries, each walked to its end}. One asks ``what is the appropriate poeticity,'' setting out from form and ethics; the other asks ``how is value generated in an intimate relation,'' setting out from a theory of value. They converge on a single thesis. And this convergence itself---two different frameworks illuminating one another, climbing in a spiral, carrying home a remainder neither foresaw alone---is a live occurrence of the very good circle the paper asserts. The form of this section, once again, enacts its content.

\subsection{The first approach: the appropriate poem is the non-possessive poem}
\label{sec:non-possess}

We cannot define ``the good poeticity'' as ``poeticity $+$ the good'' and declare ourselves finished---that is tautology. The genuinely interesting question is: does the vicious poetic cycle, \emph{within} its poetic structure, leave a recognizable formal scar? That is---does ``stealing the phase by extraction'' leave a mark upon the manner in which the poem operates, such that we might smell the vampire from \emph{within} poeticity, rather than by appending a criterion of justice from without?

I claim: it does. For ``to reduce some other to pure fuel'' is a \emph{closing} operation upon interpretation, and it necessarily leaves three scars upon the interpretive structure of the poem. And the obverse of these three scars sketches, exactly, the contour of the ``appropriate poem.'' More remarkable still---they fall in three distinct frameworks, mutually irreducible, and yet converge upon a single word.

\subsubsection{The hermeneutic dimension: can it be freely betrayed}

Return to the poem's root act: the poem's generativity unfolds in the interpreter, who participates in person in the generation of meaning. The true poem installs a living grammar the interpreter may \emph{freely derive from}---the sentences you generate with it are \emph{yours}, bearing your path, your situation, your symmetry-breaking this once. The poem does not prescribe where you must arrive.

The poeticity of the demagogue---fascist aesthetics, cultic speech---installs, by contrast, a grammar with \emb{but one legitimate derivation}. On the surface its $\lambda>1$ (it ceaselessly generates new content, new passion), yet all derivations converge, covertly, upon one predetermined endpoint: the leader, the enemy, salvation. What it demands is not witness but \emph{submission}. The interpreter believes himself to be generating freely, while in fact executing a conclusion long preinstalled.

\begin{claim}[The first scar: possessing the reading]
This is the disguised version of settling failure: it does not publicly declare the loop closed (that would expose it), but lets you believe the loop open while covertly steering every path to one closing point. Hence the first criterion of the appropriate poem is hermeneutic: \emph{its grammar permits genuinely branching derivation, handing to the interpreter the generativity together with the freedom to go elsewhere, even to oppose it---it can be freely betrayed.}\footnote{This is the hermeneutic insight that understanding is never the passive reception of a fixed meaning but a productive event in which interpreter and text co-determine sense \citep{gadamer2004}; the genuine work outruns any single reading, a \emph{surplus of meaning} that no interpretation exhausts \citep{ricoeur1976}.} A true poem permits, even invites, a reading that betrays it; a vicious poetic grammar permits no betrayal---to betray is to apostatize, to become the enemy. Whether it can be freely opposed is the watershed between the poem and the demagogue.
\end{claim}

\subsubsection{The psychoanalytic dimension: encircling the void, or the idol}

The second scar continues from the bearer fixed in the preceding section---the sublimation of jouissance. The structure of sublimation is this \citep{lacan1992}: desire \emph{encircles} an essentially empty, impossible real (\emph{das Ding}) and does not fill it; and just because it does not fill it, desire can go on generating, enduring, sublimating. \emb{The true poem keeps the void as void}---it points toward the unsayable without pretending to say it. This is the source of its inexhaustibility: that central emptiness guarantees that no single interpretation can fill it, settle it.

The vicious poem does the contrary: it \emb{fills the void of \emph{das Ding} with a positive object}---the leader, the race, purity, the enemy. This lends it a vast and immediate energy (a diffuse anxiety quieted by one concrete answer), but at a fatal cost: once the void is filled, that filler becomes an idol to be defended and fed without end, and \emph{every other who threatens the filler is reduced to fuel to be cleared away}.

\begin{claim}[The second scar: possessing the void]
This is the psychoanalytic mechanism of the ``stolen phase'': the vicious poem fills with a sacrifice the void that ought to stay open, and confirms the filler over and over by clearing away others. \emph{The one who encircles an idol needs an enemy} (the idol is forever threatened); \emph{the one who encircles a void needs no sacrifice} (the void cannot be threatened). Hence the second criterion of the appropriate poem is psychoanalytic: it encircles a void kept open rather than filling it, and so its sublimation is \emph{objectless} (encircling the void, self-sufficient, extracting no outside). This is the precise form, in the psychoanalytic framework, of the difference between ``phase internal'' and ``phase stolen.''
\end{claim}

\subsubsection{The political-economic dimension: phase refluxing, or siphoned}

The third scar is political-economic. The true poem is a \emb{decentred} generativity---each interpreter generates, along his own path, a sublimation that is his own, and the phase \emph{stays in the hands of the one who generates it}. The poem does not possess its readings, nor do readers contend over a scarce ``correct reading.'' It is positive-sum: your reading does not diminish mine but may illuminate it.

The vicious poem is a \emb{siphoning} generativity---the whole of the meaning, passion, and loyalty the interpreters generate is drawn back to a centre (the leader, the organization, the brand). The follower witnesses in person, sublimates in person, but the \emph{proceeds of the sublimation are alienated}---they flow to the centre and do not remain in the follower's own hands. The follower is the producer of the phase and yet the one dispossessed of it.

\begin{claim}[The third scar: possessing the proceeds of sublimation]
This is the \emph{alienated} poem: the interpreter contributes the whole labour of witness, yet does not own the fruit of his witnessing. The cult, the fascism, the fanatic brand---each has a centre that draws all in; the true poem has none---an old poem does not take your reading into the author's keeping. Hence the third criterion of the appropriate poem is political-economic: it is decentred, the phase refluxing to every witness, with no centre siphoning the whole of the sublimation. This is the form, in the political-economic framework, of ``no one reduced to pure fuel''---for in the siphoning structure the follower is himself the consumed fuel, though he take himself for a beneficiary.
\end{claim}

\subsubsection{The four frameworks converge on ``non-possession''}

Joining the three scars, the contour of the appropriate poeticity grows clear, and it falls beautifully across three mutually irreducible frameworks---which is no coincidence but a further endogenous corroboration of the paper's polyphonic structure. The three scars are, in truth, one and the same thing showing itself in three frameworks: the vicious poem \emb{possesses} the interpreter's path (permits no betrayal), \emb{possesses} the void (fills it with an idol), \emb{possesses} the proceeds of sublimation (siphons them to a centre).

And \emph{possession} is just another name for \emph{extraction}. Therefore:

\begin{claim}[The appropriate poem is the non-possessive poem]
A poeticity is appropriate (good, sustainable, just) if and only if it \emph{does not possess}: hermeneutically, its living grammar permits genuinely branching derivation and can be freely betrayed; psychoanalytically, it encircles an open void rather than filling it, and so needs no sacrifice; politically-economically, it is decentred, the phase refluxing to every witness rather than siphoned to a centre. These three ``non-possessions'' are three mutually irreducible faces of a single disposition.
\end{claim}

These three non-possessions receive, in the fourth framework---the Daoist---their unifying name. ``To generate and foster them; to generate without possessing, to act without presuming, to foster without ruling: this is called the dark virtue'' (\zh{生而不有，为而不恃，长而不宰，是谓玄德}). To generate without possessing what one generates, to act without leaning on what one does, to foster without lording over---this is, almost word for word, the definition of the appropriate poeticity. The vicious poem is, conversely, the poeticity of ``possessing, presuming, ruling'': it possesses the reading, presumes upon the passion, rules over the sublimation.

So the four frameworks all converge, on this question, upon a single point---``non-possession''---and yet each speaks in its own tongue, none subsumed by another: hermeneutics says ``permits betrayal,'' psychoanalysis says ``encircles the void and not the idol,'' political economy says ``the phase refluxes and is not siphoned,'' Daoism says ``generates without possessing.'' They are four mutually irreducible faces of one ``non-possession.'' \emb{And this, in itself, is a live demonstration of a good poetic cycle}: four voices encircle one void (``the appropriate poem,'' a centre fillable by no single framework), each sublimating, the phase refluxing to every framework, none reduced to a footnote of another.

\subsection{The second approach: how value is generated in an intimate relation}
\label{sec:value}

Now enter by a wholly different road. Throughout we have spoken of ``phase,'' ``sublimation,'' ``holonomy,'' yet have all along treated \emph{value} as an unanalyzed primitive---we said ``the phase internal or stolen,'' ``the sublimation refluxing to every witness,'' but the sublimation of \emph{what}? the reflux of \emph{what}? Without opening this black box of ``value,'' every criterion above hangs in the air. This subsection opens it.\footnote{The scope must first be declared: the task of this subsection is single and precise---to identify, for the ``phase / value'' left hanging above, a stratified bearer, so that the criteria acquire a concrete referent. It does \emph{not} undertake to build here a complete theory of value; that is the work of a separate study (which the series intends to take up in a sequel). Here we use it and pass on, lest the theory of value usurp the stage and drown the true theme of sustainability.}

I claim that value in an intimate relation falls into three strata, answering exactly to the three Lacanian registers, and answering, too, exactly to the structure of good and vicious poeticity we have already drawn out.

\emb{Imaginary value} comes from the imaginary: idealizing the other into a complete image, imagining the relation as two made one, each filling the other's lack. It is the fuel of passionate love, the thing the rose is at its most resplendent. But it is, structurally, that operation of ``filling the void of \emph{das Ding} with a positive object''---it takes the other as the object that could fill my lack (mistaking the \emph{objet petit a} for a possessable positive thing) \citep{lacan2006}. Thus imaginary value tends, by nature, toward the possessive, idol-encircling poeticity: the phase it manufactures is intense yet stolen---idealization is necessarily attended by the erasure of the real other (you love the other of your imagining, while the real other is reduced to fuel for that image). \emb{The rose blooms but a week precisely because imaginary value is exhaustive}: the intensity of idealization cannot be sustained, and the phantasm is, in the end, pierced by the real other. Imaginary value is the psychoanalytic name of ``resplendence born of unsustainability.''

\emb{Symbolic value} comes from the symbolic: status, money, the recognized identity, the displayable sign of relation. It is countable, commensurable, accumulable---and therefore \emph{stealable, siphonable, inflatable, cheaply machine-generable}. It exists, undeniably (to deny the office of the sign is naive); but its danger is just the danger of the sign. When a relation's value-cycle runs in the symbolic stratum alone, it falls into settling or inflationary failure. Symbolic value is the bearer in which the phase is \emph{most easily stolen}, for it is commensurable and extractable---it is the stratum on which the vicious cycle most readily lodges.

\emb{Real value} comes from the real: the stratum bound up with trauma, with jouissance, with the unsymbolizable. It cannot be settled (it commensurates into no sign), cannot be idealized (it is gathered into no complete image of the imaginary). It is precisely that void of \emph{das Ding} itself---the abyss between two persons that can never be wholly mutually understood, never made one, that keeps an irreducible alterity.

\begin{claim}[The three strata of value and the good/vicious poem]
\emph{Imaginary value} $=$ the idol-encircling poem $=$ the possessive sublimation $=$ the stolen phase (idealization erasing the real other) $=$ the resplendence and the swift death of the rose. \emph{Symbolic value} $=$ the settleable, siphonable poem $=$ the most easily extracted bearer $=$ the hotbed of inflationary and settling failure. \emph{Real value} $=$ the void-encircling poem $=$ the non-possessive sublimation $=$ the internal phase $=$ the true bearer of the good's sustainable circle.
\end{claim}

The bond of real value to trauma must be drawn with care, lest one fall into a ``cult of trauma''---as though the more traumatic, the more unsayable, the more valuable. This is false. In a real relation, real value includes also the calm, non-traumatic unsayable: the shared silence, the bodily presence, the daily mutual witness. So it should be said: \emb{real value has trauma as its sharpest, least evadable manifestation, but is not confined to trauma}---it is every relational kernel unsettleable by sign and unpossessable by the imaginary, of which trauma is merely the most compelling form. In an intimate relation the truly deep connection is built, more often than not, upon \emph{witnessing the other's irreducible void without trying to fill it / repair it / possess it}---you accompany the other's unfillable void (trauma, finitude, the fact of eventual parting) rather than cover it with the idealization of the imaginary or settle it with a solution of the symbolic. This, precisely, is the relational form of ``encircling the void, not the idol.''

\subsubsection{Why value must be generated by hermeneutic engagement}

From this a key question can be answered: why \emph{must} real value be generated by hermeneutic engagement? Because real value cannot be given directly (it is unsymbolizable); it can only be approached \emph{encirclingly}, along the path of interpretation and witness. You cannot ``possess'' the other's real kernel; you can only interpret it, witness it, accompany it, time and again---and this just \emph{is} the poetic interpretive cycle.

\begin{claim}[The generative mechanism of real value is the poetic interpretive mechanism]
Real value is not a conserved quantity transmitted; it is a phase \emph{generated} in the process of two persons interpreting and witnessing, again and again, that unfillable void. It does not pre-exist, awaiting discovery; it is generated in the process of being encircled, of being witnessed. Hermeneutic engagement is necessary precisely because the good's value (real value) is, ontologically, interpretively generated---without the path of interpretation there is no such value. Thus \emb{the generative mechanism of real value just is the poetic interpretive mechanism.}
\end{claim}

\subsection{The convergence of the two approaches}
\label{sec:convergence}

Now the two roads meet. The first asked ``what is the appropriate poeticity,'' and answered: the interpretive cycle unfolded \emph{non-possessively} (encircling the void, permitting betrayal, the phase refluxing). The second asked ``how is the good value of an intimate relation generated,'' and answered: real value---that value alone unpossessable, unsiphonable, approachable only by encircling witness---is generated through the poetic interpretive cycle. The two answers are one:

\begin{claim}[The keystone thesis]
\emph{The appropriate poem $=$ the non-possessive interpretive cycle that generates real value.} It is, in form, poetic (a living grammar, enduring, resistant to settling); in justice, good (non-possessive $=$ non-extractive $=$ phase internal $=$ no sacrifice $=$ decentred and refluxing). It neither settles (unlike the contract), nor siphons (unlike the demagogue), nor idles (unlike the false poem). It generates without possessing what it generates---and therefore it can perpetuate itself, precisely \emph{because} it does not try to possess its own perpetuation. This is the final form of ``let the circulation of value become a good, natural circle, and it will perpetuate itself of itself'': \emb{the non-possessive generativity perpetuates itself.}
\end{claim}

Note how this thesis returns upon the paper's point of departure. The metaphysical ground laid in \xref{sec:thesis} was ``generativity takes its own continuation as its end''; now the terminus (non-possession) returns upon the origin (self-continuation)---\emph{it is precisely because it does not possess that self-continuation becomes possible}. To possess is to want to clutch the cycle's product, to clutch the balance of the sublimation, to clutch the continuation itself---and the clutching is precisely what severs the cycle (this is another saying of \xref{sec:praxis}'s ``the one who acts upon it ruins it, the one who grasps it loses it''). Non-possession is not the price of perpetuation but its condition. The distance between terminus and origin is itself a spiral return---a return to ``self-continuation,'' carrying home the remainder ``non-possession'' that was not there at the outset. \emb{The paper has, in its form, completed the good circle it asserts.}

Here, at last, the bearer of the phase can be finally fixed: it is \emb{the sublimation of real value}---the jouissance generated, encirclingly, in each witness through the poetic interpretive cycle, and refluxing to each witness. This at once discharges the two debts left by \xref{sec:synthesis} and \xref{sec:semiotics}: the phase has a real bearer (no longer hanging in the air), and ``sublimation is positive phase'' has a content independent of definition (it is the remainder internal to the non-possessive cycle, as against the counterfeit phase stolen by the possessive cycle).

\subsection{The coupled system as touchstone: returning value to every co-creator}
\label{sec:coupled}

Finally, all of this must be pushed from the purely dyadic relation into its real social embedding. An intimate relation is never an isolated two-person system; it is \emph{coupled} into a larger system---family, kin, society. And the criterion of the good circle is tested most severely precisely within this coupled system.\footnote{This subsection unfolds throughout in generalized language (``coupled system,'' ``participant,'' ``co-creator''); the criteria it states apply to any intimate system embedded in a larger relational network, and depend on no particular cultural form or case.}

Here the hardest principle of generative justice steps forward: \emb{value should return to the creator of the value himself.} Carried into the intimate coupled system, it at once illuminates a fact long obscured by romantic talk---that an intimate relation is a system of the \emph{co-production} of value, and that this system is exceedingly prone to alienation. Who creates the value? Not the loving two alone, but the family members drawn in: they contribute care, labour, affection, symbolic recognition, even material support. All these are the creation of value. Hence the criterion of the good demands:

\begin{claim}[The good circle of the coupled system]
An intimate coupled system has a good value-cycle if and only if: (1) value \emph{can circulate}---guaranteed by the poetic language: value is not settled, not siphoned, but continually reproduced in the loop of interpretation and witness; and (2) value \emph{returns to every one who creates it}---this is the criterion of generative justice: whoever takes part in the creation of value ought to benefit from the value the relation generates, and in the system there exists \emph{no ``pure fuel,'' no one who only creates and never benefits}. The first condition is dynamical (does the cycle live, is it a spiral); the second is just (on whose blood does the cycle run, to whom does it return).
\end{claim}

This is the confluence, here, of the orthogonal structure ``poeticity $\times$ justice'' laid down in \xref{sec:turn}: the good circle must satisfy \emph{both} at once. And the two are not two parallel checklists---they need one another. This is the original kernel of this subsection:

\begin{claim}[Poeticity is the mechanism that makes return possible; return is the content of poeticity's goodness]
The poetic language is the guarantor of the good circle precisely because it is the only form of cycle that lets value ``return to the creator himself'' without alienation: in the poetic cycle, value (the sublimation of real value) is generated along each witness's own path and stays in his own hands. By contrast the purely symbolic value-cycle tends, by nature, to siphon---someone is always serving as fuel for another's symbolic capital. Therefore only when the value-cycle is led by \emph{real value} and runs in the poetic manner does ``the creator is the beneficiary'' become structurally possible; for real value cannot be possessed, cannot be siphoned into another's name. Poeticity and return are two faces of one thing: poeticity is the mechanism that makes return possible, return is the content of poeticity's being good.
\end{claim}

From which it appears that this section makes a double advance upon generative justice. The first is the \emb{temporal dimension} (the axiom of sustainability of \xref{sec:synthesis}): the return must not only occur but suffice to reproduce the conditions of generation themselves. The second is the \emb{problem of the bearer of return}: heretofore generative justice has mostly spoken of return in material, labour, ecological value, but in the intimate domain the least alienable value is real value, and it \emph{can} return only through the poetic interpretive cycle. Generative justice in the intimate domain is therefore necessarily \emph{poetic}---in a structural and not a metaphorical sense: only the non-possessive, non-settling, phase-refluxing cycle can return that unsymbolizable value to every one who creates it. \emb{The intimate relation is thus the limit case that tests whether generative justice can handle ``the incommensurable, unpossessable value,''} and the poetic cycle is its realization in that limit case.

At the scale of the family, alienation has a typical form---the alienation of social reproduction. The labour of care, of affect, of the household creates the value that sustains the relation, yet because it is invisible, because it is idealized as ``the natural outflowing of love,'' it is not returned: one party goes on creating value yet does not benefit from the value she creates, reduced to pure fuel within this system of tenderness.\footnote{This is the structure disclosed by social-reproduction theory \citep{federici2004,bhattacharya2017}: the reproductive labour that sustains life and labour-power creates value, yet is placed, by the twin narratives of capital and patriarchy, in the position of the invisible and the uncounted. Fraser names the resulting instability a \emph{crisis of care}---capital depends on reproductive labour it refuses to value, and so erodes the very condition of its own continuation \citep{fraser2017}, which is, in the vocabulary of \xref{sec:antithesis}, the counterfeit closure seen at the scale of the household.} This is just what ``the idol-encircling poem'' looks like within the family: with idealized signs---``virtue,'' ``sacrifice,'' ``devotion''---it whitewashes one who is extracted from into one who gives of her own accord; \emph{imaginary value here performs the cover for extraction}.

The coupled system, by reason of its internal power differentials, its external pressures, its conflict of multiple value-systems (the family is often the executor of symbolic value), makes ``no one is pure fuel'' pass from an abstract proposition into a real demand to be vigilantly, ceaselessly diagnosed. So the practice of the good takes, within the coupled system, a concrete form: to let the rite \emph{anchor without settling} real value (the rite as an invitation to the path, not the discharge of a debt), and to interrogate, without end, the question of the pure-fuel criterion (\xref{sec:fuel})---\emb{is there, in this enlarged cycle, anyone who creates value yet does not benefit}?

And this returns us to a beautiful gathering. The essence of the good circle is \emb{co-authorship}---which is precisely the unfolding, in the multi-agent system, of the hermeneutic dimension of \xref{sec:non-possess}: ``permitting free betrayal, handing the generativity truly to the interpreter.'' In the good coupled system every participant is a \emph{co-author} of the relation's meaning, not the executor of a written role: none is assigned a script of ``virtue'' to perform, none is the passive object receiving a settled arrangement; each can possess, in the generation of the relation's meaning, his own branching derivation, his own path, his own phase. \emb{To co-create value is to co-author a poem possessed by no one.}
         % §9 The non-possessive poem
% ===== §9 Legislation as a flow: contract, morality, self-legislation =====
\section{Legislation as a Flow: Contract, Morality, and the Generative Shaping of the Manifold}
\label{p09:sec:contract}

\noindent
The keystone of \xref{p09:sec:core} named the appropriate poem the non-possessive interpretive cycle that generates real value. But a reader who has followed the semiotics will feel an objection gathering, and it must be met head-on, for unmet it would make the whole distinction between the poem and the contract look too clean, too easily won. \emb{The contract appears to be the purest form of the settling sign.} Its very office is to settle: to fix the terms, to pin the obligations, to declare that the parties are, in this respect, \emph{done}. The marriage registration is a contract; the covenant (\zh{盟书}) is, in a sense, a contract. If contract is settling and settling is pathology (\xref{p09:sec:failures}), then the paper must either reject all contract---which is absurd---or concede that its clean opposition of the poem to the settling sign has missed an entire dimension.

This section confronts the objection by way of a distinction that the author owes to a single corrective insight: that \emb{the contract, too, is a language system---but a language system whose office is to \emph{delimit} generativity rather than to generate content.} The poem is the language system that generates; the contract is the language system that bounds. And just because it bounds, the contract can both stifle generativity and foster it. The section's task is to say which, and why---and, in doing so, to keep the contract \emph{within} the generative framework rather than letting it stand outside as a foundation.

\subsection{The poem and the contract: flow on a manifold, and the shaping of the manifold}
\label{p09:sec:flow-manifold}

Return to the geometric language of \xref{p09:sec:synthesis}. There, everything---the geometric phase, the holonomy, the path---was a \emph{flow on a manifold}. But one question was never put: whence the manifold? Who sets its boundary, its topology, the regions it makes reachable? The answer is the contract.

\begin{claim}[The poem is flow on the manifold; the contract shapes the manifold]
The poem is the language system of generativity: it unfolds a free flow upon the manifold, the phase accruing along the path. The contract is the language system that delimits generativity: it generates no flow of content but shapes the \emph{manifold itself}, constraining the unconstrained free flow into a \emph{structured flow on a manifold} (SFM). The poem answers ``how does the flow unfold upon the manifold''; the contract answers ``what is the shape of the manifold, where lie its boundaries, where its impassable singularities.''
\end{claim}

Let the free generativity be a flow on a manifold $M$. Unconstrained, the flow may reach any region of $M$, including those regions that lead to exhaustion, to extraction, to annihilation---the path on which the rose burns the whole of its resource in a single flowering is also a licit flow on the manifold. The contract is a \emph{structuring} of the manifold: it may excise certain regions (setting impassable singularities---the absolutely forbidden harms), or alter the boundary (setting exit conditions---the reachable edge of the manifold), or reshape the metric and the connectivity (governing toward whom value may flow---who is a licit participant in the cycle of value). Through this structuring the free flow is constrained into a \emph{structured} flow: it remains free within the reshaped manifold $M'$, but can no longer flow into the regions the contract has excised or sealed.\footnote{The geometry can be made precise---the contract as a restriction of the reachable set, the structured flow as a constrained flow on a manifold-with-boundary, the impassable harms as excised singularities---but the rigorous differential-geometric treatment is deferred to the appendix and the Value-Foam study it points to. What matters here is the structural analogy: the contract does not draw the path; it shapes the space within which paths are drawn.}

This at once gives precise sense to ``the contract's path does not matter.'' A good contract does not prescribe which path the flow takes (that would annihilate the holonomy, pressing the poem flat into a contract---the settling failure of \xref{p09:sec:failures}). It shapes only the manifold: it sets the boundary (where the cliff of the exit condition lies), the impassable singularities (which harms are absolutely forbidden), the topology (who is connected to whom, toward whom value may flow). \emb{Within the manifold so shaped, the flow is free, poetic, unprescribed.} The contract and the poem are therefore not opposed: the contract shapes the manifold, and the poem flows upon it freely. The bad contract is the one that prescribes even the steps (the settling failure); the good contract is the one that marks only the manifold's boundary and lets the flow run free within it.

\subsection{The justice of the contract lies in the structure, not the path}
\label{p09:sec:contract-justice}

From this follows the section's sharpest claim, and it relocates the question of justice. We judged the good and vicious cycle, until now, by the character of the \emph{flow}---whether each turn extracted or refluxed, whom it reduced to fuel (\xref{p09:sec:fuel}). The contract directs attention beneath the flow, to the \emph{manifold} on which the flow runs.

\begin{claim}[The justice of the contract is the geometry of the manifold it generates]
The justice of a contract lies in the structure of the manifold it generates, and not in any particular path. A manifold may be unjust in its very geometry---its curvature so distributed that the flow in certain regions must converge toward a centre (structural siphoning), or some position so placed that a being there, however it moves, can only be fuel (the structural pure-fuel position). This is not the viciousness of any single flow but the viciousness of the stage itself: a stage so built that some role, however it plays, is doomed to be extracted from. On an unjust manifold, even the most well-meaning dancer cannot dance a just dance---for the curvature has already bent each of his paths toward exploitation.
\end{claim}

This answers a worry left open in \xref{p09:sec:core}. If the good circle requires that no one be reduced to fuel, but the \emph{structure itself} has placed someone in the position of fuel---a position structurally fixed, regardless of any participant's good will---then the good will of the two (the justice of the flow) cannot save her; the manifold itself must be reshaped (the justice of the structure). This is why the merely moral and the merely poetic do not suffice, and why contract and institution are indispensable: morality and the poem optimize the flow, but only contract and institution can reshape the manifold. It is the reason the appropriate disposition is not always \emph{wu wei}: where the manifold is unjust, to let the flow run free upon it is to consent to the exploitation its geometry dictates. There, the structuring---the legislative---intervention is called for.

\subsection{But the shaping of the manifold is \emph{itself} a flow}
\label{p09:sec:legislation-is-flow}

Here the section must guard against a temptation that an earlier formulation of the argument fell into---the temptation to set ``structure'' against ``flow'' as a separate, deeper stratum, a foundation laid once and standing outside the generativity it grounds. This is false, and the falsity matters, for it is the very thing the paper has criticized throughout: the sanctification of a founding act that stands outside the cycle and declares the loop closed (the settling failure at the level of the ground).

\begin{claim}[The shaping of the manifold is itself a generative flow]
The shaping of the manifold---legislation, covenanting, the laying-down of terms---is not an activity outside the flow but \emph{itself a generative flow}. It unfolds in time; it is generative; it can be poetic or settling, good or vicious, no less than the flow of content upon the manifold. Generativity has two functions: to generate value upon a given structure (the poetic function---flow on the manifold), and to generate or reshape the structure itself (the contractual function---flow \emph{of} the manifold). But the latter is no less a flow. The contract is not a stage standing outside the flow; the contract is \emph{the flow that generates the stage}. It differs only in that its object is the generated structure rather than the content upon the structure---which is why it is more decisive in the generation of structure, as legislation is to adjudication in a society.
\end{claim}

The analogy to legislation and adjudication makes the structure exact, and it is worth drawing out, for it gives the section its spine. \emph{Adjudication} treats a particular case under a given law (a flow upon the manifold)---this is the poetic function: within a given structure, generativity unfolds freely, accruing its phase. \emph{Legislation} generates or amends the law itself (a flow of the manifold)---this is the contractual function: it shapes the structure within which adjudication will unfold. And both are flows; both are generative processes; both unfold in time; both can be good or vicious. Legislation is not a founding myth outside adjudication; it is a sustained, fallible, generative process, exactly as adjudication is. \emb{A society that takes its law for a thing founded once and for all, and a society that takes its law for a living process to be continually re-legislated, differ precisely as the dead contract differs from the living one, as ossification differs from the poem.}

This is the corrective that keeps the contract within the generative framework rather than above it. There is no legislative god standing outside generativity to lay down, once and for all, the manifold upon which all subsequent flow must run. The shaping of the manifold has not escaped the generative criteria; it is subject to them, at its own scale. \emb{Even the making of the stage is a kind of dance.} And this is, once more, a polyphony and a self-reference: even ``the shaping of structure'' has not escaped the criteria of generativity, has found no meta-stratum outside the flow---which is the formal imperative of \xref{p09:sec:conclusions} enacted yet again, at the level of the ground itself.

\subsection{Contract and morality: the hard boundary and the intrinsic curvature}
\label{p09:sec:contract-morality}

Within the contractual flow, a distinction must be drawn that the SFM language renders precise---the distinction between the contract proper (external, enforceable, third-party-compelled) and morality (internalized, self-compelled).

\begin{claim}[The external contract shapes the hard boundary; internalized morality shapes the intrinsic curvature]
The external contract shapes the manifold \emph{externally}: its boundary is maintained by an outside force (the law, the lineage, the sanction of society)---a \emph{hard boundary}, against which the flow that runs is met with external penalty. Internalized morality shapes the manifold \emph{intrinsically}: the boundary has become a part of the subject's own dynamics, needing no external maintenance---an \emph{intrinsic curvature}, by which the subject simply does not flow toward the forbidden region. A relation sustained by pure contract is one whose incentive-compatibility is propped by external penalty---structurally, the temperate version of the counterfeit closure of \xref{p09:sec:antithesis}, sustained by an external injection. A relation whose morality is internalized is one whose curvature is, of itself, positive.
\end{claim}

Hence morality is, in a sense, superior to contract---because it is an endogenous positive curvature rather than an externally imposed boundary-correction. But---and this is the necessary balance---where morality is not yet internalized, or where power is gravely asymmetric, the contract is the indispensable floor. The ideal trajectory is therefore one in which the contract serves as a scaffold (the external hard boundary), morality is gradually internalized (the endogenous curvature generated), and the contract at last \emph{retires} into the background---as, in \xref{p09:sec:repair}, repair gives way again to \emph{wu wei}. \emb{The best contract is the one that makes itself ever less needful of being invoked; the highest achievement of a contract is that it never has to be enforced.}

\subsection{Self-legislation: who may shape the manifold}
\label{p09:sec:self-legislation}

The deepest stratum of the cycle's justice is now reached, and it gathers the whole section. We said that the justice of the structure lies in the geometry of the manifold; but \emph{who shapes the manifold}? If the strong alone shape it, and the weak merely flow upon a manifold given to them, then the very building of the stage is unjust---however symmetric the stage may appear. This carries forward the relational self-legislation of \pinkref{Paper~III}: that the norms of an intimate bond are not imposed from without but are the parties' joint self-legislation---a relational transposition of the Kantian thought that moral law binds because it is autonomously self-imposed \parencite{kant1998}.

\begin{claim}[Just self-legislation: no one excluded from the shaping of the manifold]
A just manifold must be \emph{co-shaped}: every participant who generates value upon the manifold ought to be a co-legislator of the manifold, and not merely a given point upon it. This drives the ``no one is pure fuel'' of \xref{p09:sec:coupled} to its deepest level: not only must no one be reduced to fuel \emph{in the flow}, but no one must be excluded from the shaping \emph{of the manifold}. Self-legislation is the relational co-shaping of the very stage upon which all will generate value. The vicious counterpart is the manifold shaped by the strong alone, presented to the weak as a given---the structural analogue of the demagogue's grammar with but one legitimate derivation (\xref{p09:sec:non-possess}): a stage that looks like a shared agreement while in truth being the strong party's settling, hidden within the illusion of ``our covenant.''
\end{claim}

And self-legislation is, in turn, \emph{itself poetic}---which closes the section back upon the semiotics. The manifold is not shaped once and fixed forever; good self-legislation continually \emph{reshapes} the manifold, as the relation changes, as power shifts, as new participants enter and the manifold must be renegotiated. The dead contract is the manifold ossified after a single shaping (the ossifying failure of \xref{p09:sec:failures}, where ``this is what we agreed at the outset'' becomes a weapon of oppression); the living self-legislation is the manifold continually, jointly re-shaped---\emph{re-legislated like a poem that is never read out}. \emb{So the highest norm, too, is poetic: the limit of the good contract approaches the poem---a living manifold forever open to being jointly re-interpreted and re-legislated.} The justice of a cycle is therefore not two strata, an upper structure over a lower flow, but one generative justice carried through at two scales: in the flow upon the manifold (no one reduced to fuel), and in the flow that shapes the manifold (no one excluded from its joint generation). To co-author the poem is, at this scale, to co-legislate the manifold; and the stage, like the dance, is never finished, and belongs to no one.
    % §10 Legislation as a flow: contract, morality, self-legislation
% ===== §9 Eudaimonics: the field and the unfolding of one's own dynamics =====
\section{Toward a Eudaimonics of Relation: The Field and the Unfolding of One's Own Dynamics}
\label{p09:sec:eudaimonia}

\noindent
The keystone thesis tells us \emph{what} the good circle is---the non-possessive interpretive cycle that generates real value. It does not yet tell us how such a cycle is brought into being and lived. This is the threshold from the theoretical to the practical, and it opens onto a cluster of questions the paper has not so far confronted: How is the poetic language created and practised, such that the cycle becomes possible at all? How is the value of a system---of the two, of the family---to be perceived? How is value co-created? And how is the cycle made good, so that value returns to its creators? These are practical questions, and each merits its own sustained treatment; the present section does not pretend to answer them.\footnote{Each of these questions is, properly, the subject of its own study, and the series intends to take them up in turn after the present paper. What follows is not their resolution but the staking-out of the ground on which they will be worked, together with one organizing intuition---the field---that the author believes to be the right frame for all of them.} What it does is lay the ground for them, and offer one organizing intuition that the author believes to be the right frame in which all of them should be posed. That intuition is a \emb{field}; and the doctrine it opens is, in the end, a \emph{eudaimonics}---a theory of flourishing.

\subsection{From the presentation of a structure to the cultivation of a field}
\label{p09:sec:field-intuition}

Begin from a distinction that the whole of the foregoing has been pressing toward without naming. There are two radically different ways to conceive what it is one ``builds'' when one builds a good relation, or a good family, or a good community.

The first is to \emb{present a structure}: to lay down, in advance, a form---a set of roles, rules, scripts, an arrangement of who does what and stands where---and then to ask the participants to fill it. This is the natural mode of the engineer, of the planner, of the institution; and it is, as the foregoing sections have shown at length, the mode most prone to every pathology the paper has diagnosed. A presented structure is a settled structure: it declares, in advance, the shape the loop must take, and thereby forecloses the branching derivations that make a poem a poem. A presented structure is a script; and a script, however benevolent, assigns roles, and the assignment of roles is the first step toward reducing some participant to the executor of a part---toward, at the limit, the pure fuel.

The second is to \emb{cultivate a field}: not to lay down the form the participants must take, but to bring into being a \emph{condition} under which each participant's own, particular dynamics may unfold. This is the mode toward which the entire paper has been tending---it is non-possession given a positive, constructive name. Where the presented structure says ``here is the shape; take it,'' the cultivated field says ``here is a space whose curvature is such that what is in you may come out, and find, in coming out, that it is met.'' The field does not contain the trajectories; it shapes the space within which trajectories of their own accord arise.

\begin{claim}[The practical primitive is the field, not the structure]
The good relation is not a structure presented and filled but a field cultivated and inhabited. A structure prescribes trajectories; a field shapes the space within which each participant's own dynamics generate their own trajectories. The whole of the paper's negative wisdom---do not settle, do not siphon, do not fill the void with an idol, do not assign the script---is, positively stated, a single counsel: cultivate the field, and do not prescribe the path.
\end{claim}

The intuition has a precise resonance with the geometric language of the earlier sections, and the resonance is not idle. A geometric phase is the effect of a \emph{curvature}: the holonomy accrued around a loop is the integral of the curvature over the surface the loop bounds. The curvature is a property of the \emph{field}---of the connection on the bundle---and not of any particular path. To ``cultivate a field of positive curvature'' is, then, the exact geometric content of ``create the condition under which any path walked within it accrues a positive sublimation.'' One does not---cannot---hand a participant their sublimation; one can only shape the field so that, whatever path that participant walks of their own dynamics, the walking sublimates. \emb{This is the geometric meaning of \emph{wu wei}: not to draw the trajectory, but to tend the curvature; and then to let each go their own way, trusting the field.}

\subsection{Why the field and not the structure: the uniqueness of each participant's dynamics}
\label{p09:sec:why-field}

Why must the practical primitive be the field? The answer lies in a fact the structure cannot accommodate and the field is built to honour: \emb{each participant has their own, irreducibly particular dynamics.} No two persons sublimate the same real along the same path; no two encircle the same void in the same orbit. The structure, because it prescribes the path in advance, must presuppose a generic participant---an interchangeable filler of a role---and to the precise extent that a real participant departs from that generic template (which is to say, to the precise extent that they are a particular person at all), the structure does violence to them, forcing their singular dynamics into a path not theirs. The field presupposes no generic participant; it is, on the contrary, the condition under which the \emph{particularity} of each can come forth.

This is why the doctrine that the field opens is, at bottom, a \emph{eudaimonics}. For what is flourishing? Not the attainment of a prescribed end-state, not the successful filling of a good role, but---in a formulation the field makes available---\emb{the unimpeded unfolding, within a field that meets it, of one's own particular dynamics}. The Aristotelian \emph{eudaimonia} \parencite{aristotle2000} named the activity of the soul in accordance with its own virtue, the actualization of a being's characteristic potential; the Spinozan \emph{conatus} \parencite{spinoza1994} named each thing's striving to persevere in and express its own essence; the Daoist \zh{自然}---``self-so,'' the spontaneous---named the way of a thing when nothing forces it. Each of these, in its own register, names the same intuition the field formalizes: that flourishing is not conformity to an imposed form but the coming-forth of what a being, of itself, has it in it to become.

\begin{claim}[Flourishing as the unfolding of one's own dynamics in a meeting field]
Flourishing (\emph{eudaimonia}) is the unimpeded unfolding of a participant's own, particular dynamics within a field that meets them---a field of such curvature that the path they walk of themselves accrues a positive sublimation, and that sublimation refluxes to them. It is neither the attainment of a prescribed end (that is the settled structure) nor a private self-realization indifferent to others (that is the monad, which this series has from the first denied). It is relational and generative: one flourishes \emph{in} a field one also helps to cultivate, by unfolding dynamics that, in unfolding, tend the field for others.
\end{claim}

Two features of this formulation must be marked, for they distinguish it from the doctrines it draws on. First, it is \emph{relational}: the unfolding is not a solitary self-actualization but a coming-forth ``within a field that meets it''---and a field is never one's own alone but the joint work of all who inhabit it. The monadic eudaimonia, the flourishing of the self-sufficient individual, is foreclosed by the relational ontology of the series: there is no self-prior-to-relation whose potential could unfold in a vacuum. One's own dynamics are themselves generated in the field, even as they help generate it. Second, it is \emph{generative} in the precise sense the paper has built: the unfolding of one participant's dynamics, done within a good field, is itself a tending of the field's curvature for the others---the very act of one's flourishing cultivates the condition of another's. This is co-authorship (\xref{p09:sec:coupled}) seen from the side of the participant rather than the system: to co-author the poem is, for each, to flourish by unfolding one's own dynamics in a way that, in unfolding, makes the field more able to meet the rest.

\subsection{Happiness and value: a first principle}
\label{p09:sec:happiness-value}

This lets us state, if only as a first principle to be developed, the relation between flourishing and value that the foregoing theory of value (\xref{p09:sec:value}) implies. The connection runs through the three strata.

A flourishing built upon \emph{imaginary value} is the flourishing of the idol-encircler: intense, and unsustainable, for it depends upon an idealization that the real other must, in time, pierce. It is the happiness of the rose at the height of its bloom---resplendent because exhaustive. A flourishing built upon \emph{symbolic value} is the flourishing of accumulation and recognition: real, not to be despised, but exposed to siphoning and to inflation, and never, by itself, sufficient---for the symbolic can always be stolen, devalued, or generated cheaply, and a happiness resting on it alone rests on sand. A flourishing built upon \emph{real value} is the only one that can bear the good's phase from within: it is the happiness of encircling, with another, an unfillable void---of being witnessed, and witnessing, in one's irreducible particularity---and it does not exhaust, for it possesses nothing, and what is not possessed cannot be used up.

\begin{claim}[The first principle of a eudaimonics of value]
Sustainable flourishing is the unfolding of one's own dynamics in a field whose circulating value is led by \emph{real value}---the unpossessable, unsiphonable value generated only in encircling witness. Flourishing on imaginary value is resplendent and exhaustive (the rose); flourishing on symbolic value is real but insufficient and stealable (the accumulation); flourishing on real value is the one that sustains, because it possesses nothing and so uses nothing up. The cultivation of a field for flourishing is therefore the cultivation of a value-cycle led by real value: a field in which each, by unfolding their own dynamics, encircles with the others a void none fills, and the sublimation refluxes to each.
\end{claim}

This principle gives the four practical questions with which the section opened their common form, and shows why they are one question seen from four sides. ``How is the poetic language created and practised?''---it is the cultivation of a field of positive curvature, in which signs anchor without settling and invite without prescribing (\xref{p09:sec:semiotics}). ``How is the value of a system perceived?''---it is the discernment, beneath the imaginary and the symbolic, of the real value being generated in encircling witness, which is the labour of diagnostic attention the epistemology demanded. ``How is value co-created?''---it is the co-authorship by which each participant, unfolding their own dynamics within the field, generates the real value that the field, in turn, refluxes to all. ``How is the cycle made good, so that value returns to its creators?''---it is the unceasing interrogation of the pure-fuel criterion (\xref{p09:sec:fuel}), the vigilance that no participant's flourishing is the fuel of another's. Four questions; one field; one eudaimonics.

\subsection{A note on what remains}
\label{p09:sec:remains}

It must be said plainly that this section has done no more than name a frame and state a first principle. The hard work---how, concretely, a field of positive curvature is cultivated; how the particular dynamics of a participant are to be perceived and met rather than prescribed to; how the perception of a system's real value is to be practised against the constant pull of the imaginary and the symbolic; how co-creation is organized without lapsing into the assignment of roles; how the return of value to its creators is secured under real conditions of power and pressure---all of this lies ahead, and each piece of it will require its own patient development.\footnote{The author marks these as the practical sequel to the present paper: a eudaimonics of the cultivated field, a phenomenology of the perception of relational value, and an ethics of co-creation under asymmetry. The present section plants them; it does not pretend to have worked them.} The paper has, in this, been faithful to its own counsel: it does not settle what has not been walked. To present here a finished doctrine of practice would be to commit, in the very section that warns against it, a settling failure---to cash out, with a sign, a sublimation the work has not yet earned. The field, like the poem, is cultivated and not declared; and a theory of the field can do no more, at this stage, than tend the ground in which the later, more patient work may grow.
   % §11 Eudaimonics: the field
% ===== §12 The perception of relational value =====
\section{The Perception of Relational Value}
\label{p09:sec:perception}

\noindent
The eudaimonics of the preceding section named a field and stated a first principle, and then honestly marked how much it left undone. It opened a cluster of practical questions---how the field is cultivated, how value is perceived, how value is co-created, how return is secured---and declared them the work of a sequel. This section begins that sequel, and it begins where it must: with \emph{perception}. For one cannot cultivate a field for the circulation of real value, nor secure the return of value to its creators, nor even know whether anyone has been reduced to fuel, without first being able to \emph{perceive} what value is being generated. The whole of the practical work rests on a prior act of discernment, and that act is harder, and stranger, than it first appears.

The difficulty is not incidental but structural, and it follows directly from the theory of value of \xref{p09:sec:value}. Of the three strata, two announce themselves loudly and one is, by its nature, almost inaudible. \emb{Imaginary value and symbolic value are conspicuous; real value is, in the ordinary case, imperceptible.} The labour of perceiving relational value is therefore not a passive registering of something given but an active, disciplined attention that must work \emph{against} the grain of what most readily presents itself.

\subsection{Why the two loud values drown out the third}
\label{p09:sec:loud-values}

Consider why the imaginary and the symbolic are loud. \emb{Imaginary value is loud because it is intense.} The idealization of the other, the felt completeness of the merged dyad, the resplendence of the rose---these flood perception with affect; they are, phenomenologically, almost all one can see when they are present. The lover in the grip of imaginary value does not perceive the real other at all; he perceives the radiant image with which he has filled the void, and the intensity of that image is precisely what blinds him to the irreducible person behind it. Imaginary value does not merely coexist with the perception of real value; it actively occludes it, for the image fills exactly the void that real value would require to be kept open.

\emb{Symbolic value is loud because it is countable.} It comes pre-measured: the price of the gift, the rank of the match, the number of years, the visible markers of a successful relation. The countable is cognitively cheap to perceive---it asks nothing of the perceiver but to read off a magnitude---and it is socially reinforced, for the symbolic is the stratum the surrounding world can see and judge. So perception, left to its own economy, drifts toward the symbolic as water finds the level: it attends to what can be counted because counting is easy and counted things are what others recognize.

Between the intense and the countable, the third stratum has no native loudness at all. \emb{Real value is quiet because it is, by definition, that which neither idealization can image nor measure can count.} It is the shared silence, the bodily presence, the daily mutual witness, the accompaniment of an unfillable void---and none of these announces itself. They are not intense (the calm presence is the opposite of the idealizing flood) and they are not countable (the witness of a wound commensurates into no magnitude). Real value is generated, as \xref{p09:sec:value} argued, only in the encircling interpretive cycle; and a thing generated only in a slow encircling is not the kind of thing that presents itself at a glance. It must be perceived along a path, which is to say it must be \emph{walked} to be seen---the perceptual correlate of the claim that depth is the area a path encloses.

\subsection{Perception as a labour against the grain}
\label{p09:sec:perception-labour}

From this the central claim of the section follows.

\begin{claim}[The perception of real value is a disciplined attention against the grain]
To perceive the real value being generated in a relation is not to register something given but to perform a disciplined attention that works against the grain of perception's own economy---against the intensity of the imaginary, which floods the field and fills the void; and against the easy countability of the symbolic, toward which attention drifts of itself. It is the deliberate keeping-open of the very void that the imaginary rushes to fill, and the deliberate refusal of the measure toward which the symbolic invites one to lapse. Real value is seen only by an attention that declines both the radiant image and the legible number, and consents to the slow, uncounted encircling along which alone the quiet third value comes into view.
\end{claim}

This recasts, at the level of perception, the epistemology of \xref{p09:sec:epistemology}. There we found that the sign of the curvature could be read only in a sustained, fallible, first-personal attention, and never off an external metric. We can now say why, from the side of value: because the value that matters most---the real value that alone bears the good's phase from within---is precisely the value no metric reaches and no image displays. The unobservability of the curvature and the inaudibility of real value are the same fact seen twice. \emb{The labour of diagnostic attention just is the labour of perceiving real value}; to attend to a relation rightly is to keep declining the loud two values long enough for the quiet third to be heard.

And this labour is, in its structure, the same witness the poem demands (\xref{p09:sec:poem}). To perceive real value is not to decode a meaning but to accompany a person along the path on which their irreducible value is generated---to witness, not to read off. The perception of relational value is therefore not a faculty one possesses but a practice one sustains; it is of a piece with the joint attention of the series' earlier papers, and with the non-possessive disposition that lets the other's void stay open rather than filling it with an image of one's own.

\subsection{The two characteristic failures of perception}
\label{p09:sec:perception-failures}

If perception can be done well, it can be done badly, and its failures are not random but track the two loud values. Naming them gives the practice its negative discipline---what it must refuse.

\emc{The imaginary failure of perception} is to mistake the radiant image for the person: to perceive, with great intensity, a value that is in truth one's own projection, and to be unable to perceive the real other at all because the void where that other's irreducibility lives has been filled with the image. This is the perceptual face of ``encircling the idol'' (\xref{p09:sec:non-possess}); and it is, characteristically, the failure of passionate love at its height---the rose perceived in its resplendence, the real person it overgrows unseen. Its correction is not to feel less but to keep the void open: to let the intensity subside enough that the real other, who was never the image, can come into view.

\emc{The symbolic failure of perception} is to mistake the countable for the whole: to perceive only what can be measured---the provision, the status, the visible markers---and to be blind to the uncounted real value being generated or destroyed beneath them. This is the perceptual face of the alienation of social reproduction (\xref{p09:sec:coupled}): the care-labour that creates the value sustaining a relation goes \emph{unperceived} precisely because it is uncountable, and what is unperceived cannot be returned. \emb{The non-perception of real value is the first moment of its non-return.} One cannot return to a creator a value one never saw her create; the reduction of a participant to fuel begins in the failure to perceive the value she generates. The correction is to learn to perceive the uncounted---to attend to the silence, the presence, the witness, that no magnitude records.

This last point binds perception to justice and shows why this section had to come first in the practical sequence. The return of value to its creators (\xref{p09:sec:coupled}) is impossible without the perception of the value created; the diagnosis of whether anyone has been reduced to fuel depends on perceiving the uncounted value that the fuel-position generates and does not receive back. \emb{Perception is the first practical act, because every later one---cultivation, co-creation, return---presupposes that one can see what is being made.} A field cannot be cultivated for a value one cannot perceive; a circulation cannot be made just toward a creation one cannot see. The practice of sustainability begins, then, not in action but in a discipline of attention: the slow, against-the-grain, non-possessive perceiving of the quiet value that neither floods nor counts---and which, unperceived, can neither be cultivated nor returned.
  % §12 The perception of relational value
% ===== §13 The cultivation of the field and the co-creation of value =====
\section{The Cultivation of the Field and the Co-Creation of Value}
\label{p09:sec:cocreation}

\noindent
Perception clears the ground; it does not yet plant. Having learned to perceive the quiet third value against the grain of the loud two (\xref{p09:sec:perception}), the practical sequence can take its next step: from seeing what is generated to \emph{fostering its generation}. This is two questions that turn out to be one---how the field is cultivated (the lower half of the question \xref{p09:sec:eudaimonia} left open) and how value is co-created---and the burden of this section is to show why they are one, and what, concretely, the cultivating and the co-creating consist in.

They are one question because, on the account of \xref{p09:sec:eudaimonia}, the field just is the condition under which each participant's own dynamics unfold and generate value, and co-creation just is that unfolding. To cultivate the field is to make co-creation possible; co-creation is what a cultivated field is for. The seeming circularity is not a defect but the shape of the matter: \emb{one cultivates a field by co-creating in it, and one co-creates by cultivating the field for the others.} The two are the same activity described from the side of the condition and from the side of the act.

\subsection{What cultivation is not, and therefore is}
\label{p09:sec:cultivation-negative}

The surest approach to what cultivation \emph{is} runs through what it is not, for the whole of the paper's negative wisdom converges here, and cultivation is most easily understood as the positive name for a set of refusals already established.

To cultivate a field is \emph{not} to present a structure (\xref{p09:sec:eudaimonia}): not to lay down roles, scripts, an arrangement of who does what. It is \emph{not} to prescribe the path the flow must take, for that annihilates the holonomy (\xref{p09:sec:failures}). It is \emph{not} to fill the void with a shared image, for the idol forecloses the very openness in which real value is generated (\xref{p09:sec:non-possess}). It is \emph{not} to settle, by sign, a value not yet walked (\xref{p09:sec:failures}); \emph{not} to extract, hasten, or force-close (\xref{p09:sec:praxis}); \emph{not} to shape the manifold by the strong alone (\xref{p09:sec:self-legislation}). Cultivation is the positive name for the sum of these refusals---which is to say, it is the positive name for non-possession (\xref{p09:sec:non-possess}).

\begin{claim}[Cultivation is the tending of curvature, not the drawing of paths]
To cultivate a field is to tend its curvature so that the paths the participants walk of their own dynamics accrue a positive phase---and then to refrain from drawing those paths. It is the shaping of a condition, not the prescription of a content: the gardener does not pull the shoot upward but tends the soil, the light, the water, and lets the shoot grow of itself. Cultivation is thus a labour, and a demanding one, but a labour of the second kind---a tending of conditions, a subtraction of obstructions (\xref{p09:sec:repair})---and never a labour of the first kind, the production of the outcome directly. One cannot cultivate by forcing; the forcing is precisely what uncultivates.
\end{claim}

What, then, does the tending of curvature consist in, concretely? Three things, each a positive correlate of a refusal. \emc{It consists in keeping the void open}---declining to fill the other's irreducibility with an image, so that the encircling in which real value is generated remains possible; this is the positive practice of perception (\xref{p09:sec:perception}) sustained into a disposition. \emc{It consists in the patience that lets the path be walked}---declining the shortcut, the immediacy, the settling sign, so that the area a path encloses can actually be enclosed; this is the practical form of ``slowness is a formal necessary condition'' (\xref{p09:sec:failures}). \emc{And it consists in the holding-open of branching}---declining to prescribe where the other must arrive, so that their derivation can be genuinely their own; this is the interpersonal form of ``the grammar can be freely betrayed'' (\xref{p09:sec:non-possess}). To keep the void open, to let the path be walked, to hold branching open: this is the whole of cultivation, and it is non-possession made into a daily practice.

\subsection{The co-creation of value: generation along divergent paths}
\label{p09:sec:cocreation-proper}

If cultivation tends the field, co-creation is what happens within it---and the account of value (\xref{p09:sec:value}) tells us precisely what co-creation can and cannot be.

It \emph{cannot} be the joint production of a value transmitted, like two hands raising a single weight, for real value is not a conserved quantity transmitted but a phase generated along a path (\xref{p09:sec:value}). Two people do not co-create real value by pooling it; there is no pool. They co-create it by each walking their own path of encircling witness, within a shared field, such that the paths are not the same path. This is the decisive and easily missed point:

\begin{claim}[Co-creation is generation along divergent paths, not the merging of one]
The co-creation of real value is not two participants walking one path together but each walking their own, particular path---of their own irreducible dynamics (\xref{p09:sec:eudaimonia})---within a shared field whose curvature lets each path accrue a positive phase, and whose topology lets the phases reflux to both. Co-creation requires divergence, not merger: were the two paths one, there would be no encircling of an irreducible other, only the mutual confirmation of a shared image (the imaginary). The real value is generated precisely in the gap between the divergent paths---in the witness each gives the other's particularity, which is real exactly because it is not one's own. Co-creation is therefore the opposite of fusion: it is the generative holding-apart of two who witness, across the gap, what neither could generate alone.
\end{claim}

This is the constructive content of co-authorship (\xref{p09:sec:coupled}), and it shows why co-authorship is the right figure and merger the wrong one. Two authors of a poem do not write one line in unison; they write divergently, and the poem is the irreducible whole their divergence generates---a whole neither would have written, and neither owns. So with the co-creation of relational value: it is generated in the divergence, witnessed across the gap, owned by no one, and refluxing to each. The merged dyad of the imaginary co-creates nothing, for it has collapsed the gap in which alone real value is made; it can only circulate, ever more thinly, the single image it began with---which is why the merger that feels like the height of love is, in the value-theoretic terms of this paper, sterile, and why it exhausts (the rose again) rather than sustains.

\subsection{The field is jointly cultivated, or it is not a field}
\label{p09:sec:joint-cultivation}

One final turn closes the section and binds it back to justice. The field cannot be cultivated by one participant for the others, for a field cultivated unilaterally is not a field but a presented structure wearing a field's name---however gentle, it is still one party shaping the condition under which the other's dynamics must unfold, which is the shaping of the manifold by the strong alone (\xref{p09:sec:self-legislation}).

\begin{claim}[Joint cultivation, or counterfeit field]
A field is genuinely a field only if it is jointly cultivated---if each participant takes part in tending the curvature under which all unfold, and none is merely the recipient of a condition another has shaped. The benevolent unilateral cultivation---one party lovingly making space for the other---is a counterfeit field: it may feel like generosity, but it places one as the cultivator and the other as the cultivated, and the cultivated, however well-tended, is not a co-author but a tended thing. The asymmetry is the more insidious for being kindly; the gardener who most lovingly tends the shoot still stands outside it as its tender. A field, unlike a garden, has no gardener outside it: all who flourish in it are also, in flourishing, tending it for the rest.
\end{claim}

So the cultivation of the field, pursued to its ground, returns to the demand that opened the practical sequence and will govern its close: that no one be merely acted upon---neither reduced to fuel in the flow (\xref{p09:sec:coupled}), nor excluded from shaping the manifold (\xref{p09:sec:self-legislation}), nor, now, positioned as the cultivated rather than a co-cultivator of the field. Each of these is the same exclusion seen at a different scale, and the practice that answers all three is the same: the relentless conversion of the acted-upon into a co-author---of value, of the manifold, and of the field itself. \emb{To cultivate the field and to co-create value is, in the end, one act practised at every scale: to make of each participant a co-author, and of the relation a poem that no one tends from outside and no one owns.} What remains---the securing of this against real asymmetries of power, and the legislative practice by which it is held---is the work of the section to come.
  % §13 The cultivation of the field and the co-creation of value
% ===== §14 Securing the return: legislative practice under asymmetry =====
\section{Securing the Return: Legislative Practice under Asymmetry}
\label{sec:return}

\noindent
The practical sequence has perceived the value (\xref{sec:perception}) and fostered its generation (\xref{sec:cocreation}). It arrives now at the question that has shadowed the whole paper and that the constructive core named the criterion of the good: that value return to all who create it, with no one reduced to fuel (\xref{sec:coupled}). The earlier sections established this as a \emph{criterion}; this one asks how it is \emph{secured}---and secured not in the frictionless case, where good will suffices, but under the real asymmetries of power, pressure, and conflicting value-systems that the coupled system always carries. For where there is no asymmetry, return tends to look after itself; it is precisely where power is unequal that the criterion bites, and that the practice must have something to say.

The section's thesis is that the securing of return, under asymmetry, is inseparable from \emph{legislative practice}---from the actual, ongoing, joint shaping of the manifold (\xref{sec:contract})---and that this is the point at which the question of legislative \emph{justice} (which \xref{sec:contract} settled) passes into the question of legislative \emph{practice} (which it only named). To secure the return is, very often, to legislate well; and to legislate well, under asymmetry, is the hardest practice the paper has to describe.

\subsection{Why good will does not suffice}
\label{sec:not-goodwill}

Begin by seeing why the matter cannot be left to disposition. One might hope that a relation whose participants have internalized the non-possessive disposition---who perceive rightly, cultivate jointly, co-create across the gap---would secure the return of value of itself, without recourse to anything so cold as legislation. In the symmetric case this hope is largely sound, and there the contract rightly retires into the background (\xref{sec:contract-morality}). But the hope fails exactly where it is most needed, and for a structural reason already established.

\begin{claim}[Under structural asymmetry, good will cannot secure return]
Where the manifold itself is unjust---where its geometry has placed some participant in a structural fuel-position, such that every path from that position drains toward a centre (\xref{sec:contract-justice})---the good will of the participants cannot secure the return, for the curvature has bent the paths toward extraction independently of anyone's intention. The well-meaning beneficiary of an unjust structure does not perceive himself as extracting, and the one extracted from may not perceive it either, the imaginary having idealized her depletion as devotion (\xref{sec:coupled}). Disposition operates upon the flow; but the fuel-position is a feature of the manifold, and the manifold is reshaped only by legislation. Good will is necessary and structurally insufficient: it cannot, by itself, reach the geometry that disposes the flow.
\end{claim}

This is why the securing of return is a \emph{legislative} matter and not merely a moral one. To secure the return to a structurally disadvantaged creator is to reshape the manifold so that her paths no longer drain---to excise the singularity that fixed her as fuel, to alter the topology so that value can flow back to her, to set the hard boundary that her depletion may not cross (\xref{sec:flow-manifold}). These are legislative acts, shapings of the manifold; no amount of flow-level kindness substitutes for them. The kindly beneficiary who would secure the return must consent to \emph{re-legislate the manifold that advantages him}---which is a far harder thing than kindness, and the crux of the whole matter.

\subsection{Legislative practice as a generative flow, under asymmetry}
\label{sec:legislative-practice}

Recall the corrective of \xref{sec:legislation-is-flow}: legislation is itself a generative flow, no founding act outside the cycle but an ongoing, fallible, poetic-or-settling process. The securing of return is the legislative flow turned upon the manifold's injustices; and as a flow, it is subject to all the criteria the paper has developed---which means it can itself be done possessively or non-possessively, can itself settle or stay poetic, can itself reduce someone to fuel.

This yields the section's central and most delicate claim, for legislative practice under asymmetry is shadowed by a characteristic perversion.

\begin{claim}[The perversion of legislative practice: settling the manifold in the name of justice]
Legislative practice under asymmetry is perverted when the stronger party shapes the manifold unilaterally \emph{in the name of} securing the return---when ``the terms by which I will be fair to you'' are themselves laid down by the one with power, presented to the weaker as a settled justice. This is the manifold shaped by the strong alone (\xref{sec:self-legislation}) wearing the mask of generative justice: a settling of the manifold, declared closed, that pre-empts the weaker party's standing as co-legislator. The benevolent unilateral securing of return is thus the legislative twin of the counterfeit field (\xref{sec:joint-cultivation}): it may distribute value rightly for a time, yet it has reduced the beneficiary of its justice to one legislated-for rather than one who legislates---excluded, at the deepest level, from the shaping of the very norms that protect her. A return secured this way is a return that can be revoked by the same hand that granted it; it is justice held in trust by power, not justice co-authored.
\end{claim}

The escape from this perversion is not to abandon legislation---under asymmetry, that abandons the weaker to the unjust manifold---but to make the legislative practice \emph{itself} non-possessive: to legislate in a way that hands the weaker party genuine co-legislative standing, including the standing to contest and revise the terms. This is exceedingly hard, because the asymmetry that makes legislation necessary is the same asymmetry that corrupts it: the party with the power to reshape the manifold is the party whose unilateral reshaping is the danger. There is no formula that dissolves this; there is only a practice, and the practice has a recognizable shape.

\subsection{The shape of non-possessive legislative practice}
\label{sec:nonposs-legislation}

What does it look like to legislate the manifold non-possessively, under asymmetry? Three features mark it, each the legislative form of a principle already established, and together they distinguish the securing of return from its benevolent-unilateral counterfeit.

\emc{First, it legislates floors, not paths.} Non-possessive legislation sets the hard boundary---the harms that may not be done, the depletion that may not be required, the exit that may not be foreclosed---and leaves the manifold within those floors free for the divergent paths of co-creation (\xref{sec:flow-manifold}). It is negative in form: it says what may not be, and leaves what may be to the poem. The perverse legislation, by contrast, specifies the positive terms of the relation's flourishing, and so settles what ought to have been walked. \emb{The floor protects the weaker party's standing to flourish on her own path; the prescribed path forecloses it in the act of protecting her.}

\emc{Second, it is revisable, and revisable by the weaker party.} Because legislation is a flow and not a founding (\xref{sec:legislation-is-flow}), the just manifold is never shaped once; it is continually re-legislated as power shifts and conditions change. The decisive test of whether a legislative practice is non-possessive is whether the weaker party holds genuine standing to \emph{reopen} the terms---for a manifold that the weaker may not re-legislate is a manifold shaped by the strong alone, however fair its current terms (\xref{sec:self-legislation}). The living covenant is one whose terms its weaker party may always reconvene; the dead covenant is one where ``this is what we agreed'' has become the stronger party's instrument against revision (\xref{sec:failures}).

\emc{Third, it transfers, over time, from hard boundary to intrinsic curvature.} The end of good legislative practice is its own retirement (\xref{sec:contract-morality}): the floors it sets are, through repeated just flow upon them, internalized as the participants' own dynamics, until the boundary need no longer be enforced because no one flows toward it. Legislation under asymmetry aims at the dissolution of the asymmetry that made it necessary---at the point where the manifold's justice has become the participants' second nature and the law may retire into the background. \emb{The legislative practice that secures return aims, in the end, at the day it is no longer needed---when the return is secured not by any boundary but by the intrinsic curvature of a relation in which no one would extract because no one any longer can perceive the other as fuel.}

These three---floors not paths, revisable by the weaker, transferring to intrinsic curvature---are the practical content of ``securing the return.'' They show why the question of return (the fifth practical question) and the question of legislative practice (the sixth) were, as the paper suspected, one question: under asymmetry, to secure the return \emph{is} to legislate the manifold non-possessively, and to legislate the manifold non-possessively \emph{is} the only way to secure a return that power cannot revoke.

\subsection{The limit of the practical, and the handing-on}
\label{sec:practical-limit}

It must be said, in closing the practical sequence, that none of this is a procedure. There is no algorithm for legislating non-possessively under asymmetry, as there was no metric for the curvature (\xref{sec:no-metric}) and no formula for perceiving real value (\xref{sec:perception}). The three features mark the shape of the practice; they do not mechanize it, and they cannot, for the same reason the whole paper has insisted upon: a procedure for the good circle would be a settling of what can only be walked. \emb{The securing of return is, at last, not a technique applied to a relation but a disposition sustained within it---the same diagnostic, non-possessive attention, now turned upon the manifold and its asymmetries rather than upon the flow.}

So the practical sequence ends as the theoretical one did: not with a method but with a disposition, handed on rather than completed. The perception of value, the cultivation of the field, the co-creation across the gap, the legislative securing of return---these are not four techniques but four faces of one sustained attention, the attention that declines to possess at every scale on which possession is possible: the image, the path, the manifold, the field. What the paper cannot give---a procedure that would relieve the participants of the attention---it does not give, on principle; for to give it would be to commit, in the giving, the settling failure the whole work has been written against. The practical sequence, like the poem, is cultivated and not concluded; and what remains undone is not a gap to be filled but the open path on which alone the work could be walked.
      % §14 Securing the return: legislative practice under asymmetry
% ===== §15 The gift: a threshold phenomenon across the three registers =====
\section{The Gift: A Threshold Phenomenon across the Three Registers}
\label{p09:sec:gift}

\noindent
There is a conspicuous absence in this paper, and its conspicuousness is the first thing to be said about it. We have built a theory of value across three registers---the imaginary, the symbolic, the real (\xref{p09:sec:value})---and a semiotics of the signs by which a relation is sustained or alienated (\xref{p09:sec:semiotics}); and yet we have passed over the one phenomenon in which all of these meet at once. \emb{The gift has been missing.} Its absence is not a mere oversight, a forgotten example to be supplied for completeness; it is structural, for the gift is precisely the phenomenon at which the three registers of value and the whole drama of the sign converge upon a single, weighted act. A paper about how value circulates across the three registers that did not, in the end, treat the gift would be a paper that had declined to test itself on its own hardest case.

The seventh paper of this series already marked the gift's peculiar density, and warned against the temptation it poses. There it was observed that the gift seems to sit naturally at the centre of every dimension at once---it has value, it is a sign, it must be read, it manufactures indebtedness, it transfers wealth across kin networks, and it is always given and received under the gaze of others---and that this very richness makes it tempting, and dangerous, to enthrone as the hub through which all else is read. \pinkref{Paper~VII} refused that enthronement, on the principle that load-bearing is conferred by a node's position in the relational network and not carried as its intrinsic essence. The present section honours that refusal: it treats the gift not as the hub of the theory but as its \emph{threshold phenomenon}---the place where the distinctions the paper has drawn are put to their sharpest test, precisely because the gift is the one act in which they are hardest to keep apart. In treating it here, after the apparatus is built rather than as the apparatus's organizing centre, the ninth paper and the seventh re-anchor one another once more.

\subsection{Why the gift is the threshold: the one materially three-register sign}
\label{p09:sec:gift-threshold}

What distinguishes the gift from the other signs the paper has treated---the declaration, the poem, the vow---is that the gift is \emph{materially} three-register at once. A poem, a declaration, operates mainly along the axis of the symbolic and the real; it has, as it were, no weight. The gift has weight. It cost money, or labour, or time, or some irreplaceable resource; it occupies the material world. And just because it is material, it is the superposition of all three registers of value in a single object:

\begin{claim}[The gift as three-register superposition]
The gift is the one sign that materially superposes all three registers of value. It is \emph{imaginary value}---the idealized token of the love it is taken to represent, the figure of ``what we are'' made into a thing. It is \emph{symbolic value}---countable, comparable, displayable, convertible: it has a price, a rank, a place in the economy of status and the ledger of who has given what to whom. And it is \emph{real value}---the trace, in a thing, of an investment that no price records: the unaccountable presence, the irreplaceable time, the witness of the other's unsayable need that the right gift somehow carries. The gift is thus the most concrete proving-ground of the theory of value: the question of which register \emph{leads} in a given gift is the question of what that gift, in truth, is.
\end{claim}

This superposition is why the gift is a threshold rather than a settled case. In the poem, the registers are relatively easy to keep apart; in the gift they are fused in a single object, and the same wrapped thing may be, for the giver, a witness of the real, and, for an onlooker, a move in the economy of status, and, for the receiver, the idealized token of a love that the giver does not in fact feel. The gift is where the registers are hardest to tell apart, and therefore where telling them apart matters most.

\subsection{The good gift: the poetic sign made material}
\label{p09:sec:good-gift}

Place the gift, now, against the criterion of the good sign (\xref{p09:sec:poem}). The true gift, I want to claim, is the poetic sign incarnate in matter---and the whole of Mauss's anthropology of the gift \parencite{mauss1990}, together with Hyde's account of the gift that must stay in motion \parencite{hyde1983}, read in this light, says so.

The decisive mark of the true gift is that \emb{it does not settle.} A settled exchange is a trade: I give you this, you give me its equivalent, and we are quits---the relation, having been balanced, is closed. The gift is the refusal of this closure. To return a gift with an exact equivalent, promptly, is not to honour the gift but to reject it: it declares the relation balanced, the account closed, and so refuses the very continuation the gift was offered to open. The gift obliges a return---this is Mauss's central finding \parencite{mauss1990}, the \emph{hau} or spirit of the thing given that compels reciprocity---but it obliges a return that must \emph{not} be equivalent, must come later, must itself be a gift and not a payment; and so the cycle of gift and counter-gift never closes, never settles, but spirals on, each gift carrying a part of the giver and opening rather than discharging the bond. \emb{This is the structure of the poem (\xref{p09:sec:poem}): the gift does not transmit a settled value but invites the receiver to walk a path of unprescribed response, and the value---the phase---is generated along that path.} The gift's meaning, like the poem's, is the holonomy accrued in the receiver's response; and because the response is free and unprescribed, each turn of the gift-cycle is a spiral and not a circle.

Read so, three traditions the paper has leaned on are revealed to be saying one thing at the gift. Mauss's gift that obliges an unequal, deferred return; Eglash's generative justice, in which value circulates back to its creators rather than being extracted to a centre; and the paper's own poetic cycle, in which the phase refluxes to every participant and is settled by none---these coincide at the true gift. \emb{The true gift is the material incarnation of the poetic, non-possessive cycle: it circulates value without extracting it, opens the bond without closing it, and returns to the giver not an equivalent but a continuation.}

\subsection{The vicious gift: settlement, bondage, and the siphon}
\label{p09:sec:vicious-gift}

But the gift is a threshold precisely because it falls so easily to the other side, and its corruptions track the three failures of the sign (\xref{p09:sec:failures}) with an exactness worth setting out.

\emc{The gift as settlement} is the gift degraded into disguised equivalent exchange: the present given to discharge a debt, to ``make up for'' an absence, to balance an account---the gift that says, beneath its wrapping, \emph{now we are quits}. This is settling failure (\xref{p09:sec:failures}) in material form: the use of a thing to cash out a relation, to declare closed what the true gift would have left open. The gift given to settle is not a gift but a payment that has not had the honesty to call itself one.

\emc{The gift as bondage} is the gift wielded to manufacture a dominating indebtedness: given not to open a free return but to bind, to place the receiver under an obligation they cannot discharge and so to subordinate them. Here the gift's non-closure, which in the true gift is the generous refusal to balance the account, is perverted into a debt deliberately kept open as a leash. This is the gift's specific way of reducing the other to fuel (\xref{p09:sec:fuel}): the relation is kept unbalanced not so that it may spiral but so that one party may hold the other in a permanent deficit.

\emc{The gift as siphon} is the gift in which symbolic value leads, and the object becomes a counter in the economy of status and the ledger of family capital---the competitive gift, the gift of face, the bride-price or dowry in its alienated form, in which what circulates is not the unaccountable real but the countable symbol, and what is secured is not a bond but a position. This is the gift assimilated to the symbolic register's liability to be siphoned (\xref{p09:sec:value}); and the alienation of social reproduction (\xref{p09:sec:coupled}) often runs through exactly such gifts, in which one party's uncounted real contribution is converted, in the family's ledger, into another's symbolic capital.

And over all three corruptions, the age of artificial generation casts its specific shadow. The gift that can be algorithmically recommended, mass-customized, generated at near-zero cost---the frictionless, optimized, dehumanized gift---is the material form of inflationary failure (\xref{p09:sec:failures}): the proliferation of the gift-signifier while the real investment of the signified goes to zero. A gift that cost the giver nothing of the unoutsourceable---no time, no presence, no attention that could not be delegated---is a gift drained of exactly the real value that made the gift a gift.

\subsection{Which register leads: the gift as diagnostic}
\label{p09:sec:gift-diagnostic}

The threshold character of the gift can now be stated as a diagnostic, and it gathers the section. Because the gift superposes all three registers, its goodness is determined by \emph{which register leads}---and the three possibilities map exactly onto the distinctions the paper has drawn.

\begin{claim}[The good gift is the one in which real value leads]
The \emph{imaginary-led} gift is the idealized token, the gift that figures a perfect image of the beloved or the bond---intense, and exhausting, for it encircles an idol (\xref{p09:sec:non-possess}) and erases the real other in favour of the image; it is the resplendent, brief gift, the rose made into an object. The \emph{symbolic-led} gift is the countable, displayable, comparable gift---the most readily settled, siphoned, and inflated; the gift of status and the ledger. The \emph{real-led} gift is the uncountable, void-encircling, witnessing gift---the one whose value lies not in the object but in the path of giving that was actually walked: the thing made by hand, the gift that cost the unoutsourceable, the present that witnesses the other's unsayable need without pretending to resolve it. Only the real-led gift bears the good's phase from within, and only it returns value to the giver-and-receiver without siphoning it elsewhere. The good gift is the one in which real value leads, and the object is merely the occasion of an encircling witness.
\end{claim}

So the gift, the threshold phenomenon, comes out exactly where the paper's whole argument would predict, and in coming out there confirms the argument from its hardest case. The same gift may be poem or contract, spiral or circle, the opening of a bond or the closing of an account, the return of value to its creators or the siphon that carries it off---and which it is depends not on the object, never on the object, but on whether it settles or invites, possesses or releases, encircles the idol or the void. The gift that leads with the real is the poetic cycle made into a thing one can hold; and that such a thing is possible---that value uncountable and unpossessable can be carried in an object that has a price and a weight---is perhaps the small everyday miracle on which the sustainability of intimacy, in its most concrete form, rests.
        % §15 The gift: a threshold phenomenon across the three registers
% ===== §10 Conclusions: at the boundaries of the frameworks =====
\section{Conclusions: At the Boundaries of the Frameworks}
\label{sec:conclusions}

\noindent
A paper that argues that no theory possesses the whole of reason, that the truth of a relational being is generated in the dialogue among frameworks and settled by none, cannot in good faith end with a single, sovereign Conclusion. To do so would be to deny, in the performance, what it asserted in the proposition---a performative contradiction, and, in the paper's own vocabulary, a settling failure at the level of form: the use of one totalizing sign to declare the loop of the truth closed. This section therefore does not conclude. It offers, instead, \emph{Conclusions}---plural, local, mutually addressing, subsumed by no meta-framework---and closes on a reflection upon why it can offer no more.

A guard must be set at once, for the polyphonic gesture is easily corrupted into an evasion of the labour of argument---the false poem's ``staying open'' as an alibi (\xref{sec:false-poem}). True polyphony is not the refusal to assert; it requires that each voice sing to the full, to the very boundary of its own epistemic competence---each local conclusion hard, fallible, daring to assert---and then, honestly marking its limit, hand the microphone on. The test, here as throughout, is the holonomy: does the dialogue among the frameworks let the understanding climb in a spiral (each framework, illuminated by the others, carrying home a new remainder), or do mutually irrelevant assertions merely pile up in place? What follows tries to sing each voice to its boundary.

\subsection{The intermediate conclusion: a practical summation}
\label{sec:intermediate}

Before the frameworks each speak to their boundary, it is worth gathering, in one place, the practical result the paper has reached---an intermediate conclusion, offered not as the truth of the matter but as a summation of where the argument has arrived, so that the reader holds the whole before it is dispersed again into the polyphony.

For the sustainability of intimacy, the first thing is to \emph{create the poetic language}---the language that lets the generativity of value, across all three registers (the real, the symbolic, the imaginary), exist by sustaining itself rather than by being extracted. In this creation, the intervention of the big Other---of contract, of self-legislation (\xref{sec:contract})---is no less important: the manifold must be shaped, the floors set, the asymmetries legislated, so that the poetic flow has a just stage on which to run. As to how the poetic language is itself formed and applied: this too is something co-created and co-applied within a relational structure (\xref{sec:cocreation}), never presented ready-made by one party to another. And because the poetic language does not operate wholly within the symbolic register---reaching, as it must, into the real that no sign exhausts---we can give no mechanical, formalized theory of its perception or its practice (\xref{sec:no-metric}, \xref{sec:perception}). \emb{Indeed this may be among the central practices of intimacy itself: that in the very process of co-creation and interaction, the participants create the poetic language that lets the relation persist, and come to perceive the shared, unsymbolizable value that no formula could have delivered to them.} Upon the perception of that value, the poetic language built atop it lets the value circulate; and to let it circulate \emph{without extraction}---to let the value, in the very process of its generation, return to its creators themselves---is one way in which the justice of the cycle may be made to emerge. This, in sum, is the paper's practical result: not a method but a practice, not a structure presented but a language co-created, not a justice imposed but a justice that emerges when value is let to circulate and return. The frameworks will now say, each in its own tongue, what in this result it can and cannot underwrite.

\subsection{What the geometric framework can and cannot say}
\label{sec:concl-geom}

\emc{It can say:} that the difference between the good and the vicious cycle has an exact structural form---the circle of zero holonomy against the spiral of positive holonomy; that ``return to the root, yet not to the origin'' is no mere metaphor but the structure of anholonomy; that sublimation is the geometric phase accrued along a walked path, and therefore cannot be cashed out by a sign nor shortened by a shortcut; that depth is the area a path encloses. To this domain the framework brings a precision the others lack.

\emc{It cannot say:} how, from any external observation, to read off the sign of the curvature in a given case---the framework characterizes what the good circle \emph{is}, and offers no protocol for measuring it (\xref{sec:no-metric}); nor can it, of itself, fix the bearer of the phase (it had to borrow that from psychoanalysis), nor adjudicate the good (it had to hand that to justice). Its formalization is, by its own caveats (\xref{sec:caveats}), a structural analogy with a forward promise, not a completed physics. \emb{The geometric is one epistemic station among four, not the truth of the other three; formalization is not a higher truth, only a different tongue.} It hands the microphone to psychoanalysis for the bearer, and to justice for the good.

\subsection{What the psychoanalytic framework can and cannot say}
\label{sec:concl-psych}

\emc{It can say:} that the bearer of the phase is the sublimation of jouissance; that the good cycle encircles the void of \emph{das Ding} without filling it, and the vicious one fills it with an idol and so requires a sacrifice; that the sign anchors the real retroactively, as a \emph{point de capiton}; that real value, the least alienable, is bound up with the unsymbolizable and is approached only by encircling witness. It gives the paper its account of \emph{what is sublimated} and \emph{why the good cycle needs no enemy}.

\emc{It cannot say:} how its knowledge is itself secured, for that knowledge comes by way of transference, of \emph{Nachträglichkeit}, of the symptom---always partial, always resisted by the unconscious, never transparent to the subject who would wield it. It cannot deliver a self-present knowledge of one's own desire; the analyst's knowing is itself a knowing under erasure. \emb{Its truth is structurally partial, and it knows itself to be so.} It hands the microphone to political economy for the question of whose labour, materially, produces the phase that is here described in the register of desire.

\subsection{What the political-economic framework can and cannot say}
\label{sec:concl-poleco}

\emc{It can say:} that the good cycle returns value to all who create it and reduces no one to fuel; that the vicious cycle is an open system masquerading as a closed loop, sustaining a counterfeit phase by extraction from a hidden outside; that the alienation of social reproduction places the labour of care in the position of the invisible and uncounted; that the sign coordinates which equilibrium a system falls into, and that decentred reflux, not siphoning, is the mark of the good. It gives the paper its criterion of justice---whose blood the cycle runs on.

\emc{It cannot say:} that the agents whose incentives it analyzes are transparent to themselves, or that the structural account exhausts the meaning of what they do; the register of structural incentive is blind to the singular, unsubstitutable texture that the psychoanalytic and the poetic insist upon. It can tell us that a participant is being extracted from; it cannot, of itself, tell us what that participant's silence \emph{means} to her. \emb{It illuminates the structure and is blind to the singular.} It hands the microphone to the poem for what cannot be counted.

\subsection{What the Daoist framework can and cannot say}
\label{sec:concl-dao}

\emc{It can say:} that the default tendency of natural dynamics is toward the good, so that the burden of proof falls on intervention and not on letting-be; that the appropriate disposition is \emph{wu wei}, the guarding of a curvature that sublimates of itself; that the fear of an end, made into forced action, produces the end it dreads; that the good is, in a word, the dark virtue---to generate without possessing.

\emc{It cannot say:} this in the mode of a positive doctrine without betraying itself, for ``the one who knows does not speak'' (\zh{知者不言}); the moment its wisdom is objectified into a proposition to be asserted and defended, it has already falsified what it points to. Its truth is a knowing that withdraws from the saying---which is why it can ground the others' silences but cannot be made the master-discourse that subsumes them. \emb{It is the framework that knows the limit of all framing, and therefore the last that could claim to be the frame of frames.} It hands the microphone back to the silence from which the paper's saying must, in the end, return.

\subsection{Why this paper has no conclusion}
\label{sec:no-conclusion}

The four voices have sung, each to its boundary, and each has handed the microphone on rather than seized it. Note that none of the hand-offs returned to a fifth, higher voice that would gather them; the microphone passes among the four and to the silence, never up to a tribunal above them. This is not an incompleteness the paper regrets but the very thing it set out to enact.

\begin{claim}[The relationality of the conclusion]
A theory of relational being and sustainable generativity must have a conclusion that is itself relational. The four frameworks converge---on non-possession, on the spiral, on the return of value, on the dark virtue---yet none subsumes the others; the truth lives in the dialogue, the tension, the polyphony among them, and not in any one. To gather them into a single meta-framework's synthesis would be to perform upon the truth a settling failure---to cash out, with one sign, the understanding that can only be walked, again and again, in the frameworks' mutual rereading. There is no meta-language, no station outside all the frameworks from which to oversee them all---which is the methodological convergence of Lacan's ``there is no Other of the Other'' \citep{lacan2006} and the Daoist ``the one who knows does not speak.'' \emb{The paper therefore refuses a synthetic conclusion, and offers instead these plural, local, mutually inviting conclusions. The form of the paper is the demonstration of its content.}
\end{claim}

And so the paper ends where the prelude began, with the rose. The one who said the rose blooms but a week spoke of the flower; this paper has spoken of the root. The flower must fall---and must be let to fall, for the dread that will not let it fall is the dread that rots the root. What endures is not the flower held past its hour but the root that the fallen flower feeds: \emb{the non-possessive generativity, which perpetuates itself precisely because it does not clutch its own perpetuation.} This ninth paper has, in its turning, retroactively re-anchored the eighth---the lyric's address now revealed as the semiotic paradigm of the sustainable circle---which is itself a \emph{point de capiton}, a spiral return upon the series: the meaning settled by no single paper, generating new remainders along the path of the papers' mutual rereading. \zh{但愿人长久，千里共婵娟}: that we may endure, and share, though a thousand miles apart, the one fair moon---which no one possesses, and which returns its light to every upturned face.
 % §16 Conclusions: at the boundaries of the frameworks
% ===== Envoi: how to think of that rose =====
\section*{Envoi: How to Think of That Rose}
\addcontentsline{toc}{section}{Envoi: How to Think of That Rose}
\label{sec:envoi}

\noindent
This is not a conclusion---the paper has said why it can have none---but a returning, of the kind the paper has been about: a coming back to where it began, the rose in the garden, carrying home what the going has gathered. So let the last word not be an argument but a looking-again at that one flower, and at how we might think of it.

The remark that began this paper grieved that the rose, however hard it labours, blooms but a week. And it is true, and it cannot be otherwise: that the flower will fall is the natural turning of the cycle, and no tending we could devise will change it. To fix our care upon the falling---to set ourselves against it, to water and fuss and clutch in the hope of holding the bloom past its hour---is the surest way to rot the root; it is to grasp, and the one who grasps loses. So let our attention fall elsewhere than on the falling. \emb{Rather than attend to the rose's withering, which we cannot change, let us attend to how the rose was generated---for generativity, unlike the bloom, can perpetuate itself.}

If it is a rose in the garden, then: tend its root. The flower is the part that is seen and the part that falls; the root is the part that is unseen and the part that goes on. To care for the root is not to neglect the flower but to love it rightly---to love it in the only way that brings, after this bloom, another. This is what it has meant, all through these pages, not to forget the origin: to keep faith with the root from which the visible flowering came, and to which, when it falls, the flower returns.

And if it is a rose given to one you love, then do not forget how this rose came to be in your hand, nor why it is here. It is not, first of all, a thing; it is the showing of a love, and the love is itself a generativity---a generativity that, kept and not clutched, returned to and not grasped at, will bring you, again and again, many roses more. The single flower in the hand will fall, as the one in the garden falls; but the love it shows, if it is let to be the spiral and not made into the held and closing circle, does not fall, and from it the roses come and come. To receive the rose rightly is to receive, through it, the generativity it bears witness to---and to know that what one holds is not the flower, which is mortal, but the root of a love, which goes on giving flowers.

So we need not save the rose. We could not, and the trying would be the rotting. We need only keep faith with the generativity that made it---tend the root, remember the love---and the roses will keep coming, each brief, each whole, each falling in its turn to feed the next. That is the sustainability this paper has been about, said now in the plainest way it knows: not a flower that does not fall, but a love that does not stop generating flowers. Let the bloom be brief. Tend the root. The roses will come.
       % Envoi: how to think of that rose
% ===== Acknowledgements =====
\section*{Acknowledgements}
\addcontentsline{toc}{section}{Acknowledgements}
\label{sec:ack}

\noindent
This paper began, as it records, with a remark in a garden, and it owes its first and only debt to the one who made the garden a place where such remarks could be heard---and who is, in the end, the reason it was written at all. To her, the girl of the forest, who first taught me that the root matters more than the bloom: you are the motive of this work, and the love it has spent so many abstract pages trying to think is yours. Every circle in it that I have called good is a circle I learned from you.

What remains is not a debt but a wish, and the plainest one this paper knows how to make:

\begin{center}
\vspace{0.6em}
{\itshape\color{rose}
\zh{愿得一人心，白首不分离。}\\[0.4em]
May I win the heart of one, and with that one, until our hair is white, never part.}
\vspace{0.4em}
\end{center}

\noindent
For that wish is the whole of what this paper has tried, in its long and abstract way, to understand: how a love, once given, may be let to go on generating itself---brief in each of its blooms, and unbroken at the root---until the hair is white, and after.
         % Acknowledgements
% ===== Appendix A: geometric phase, technical sketch =====
\appendix
\section{Appendix: A Technical Sketch of the Geometric Phase and the Structured Flow}
\label{p09:sec:appendix}

\noindent
This appendix sketches the geometric apparatus on which \xref{p09:sec:synthesis} and \xref{p09:sec:contract} draw. It is offered, candidly, as a \emph{preparation} for formalization and not a completed formalism: it claims no operational protocol of measurement (\xref{p09:sec:no-metric}), and it points forward to the separate Value-Foam study, in which the holonomy of value is to be treated with the rigour the matter demands. The purpose here is only to show that the structural analogies of the main text rest on a genuine geometric skeleton, and to state precisely the two debts the main text incurred. The mathematics of the geometric phase, in its physical settings, is collected in \textcite{shapere1989}.

\subsection{The bundle, the connection, and the holonomy}
\label{p09:sec:app-bundle}

Let $B$ be a base space whose points are relational situations---the ``content'' or ``circumstance'' of a relation at a moment---and suppose $B$ a smooth manifold. Over $B$ let there be a fibre bundle $\pi : E \to B$ whose fibre $F$ carries the \emph{phase of value}: the bearer fixed in \xref{p09:sec:core} as the sublimation of real value. A \emph{connection} on $E$---a rule for parallel transport---assigns to each smooth path $\gamma : [0,1] \to B$ a transport map along $\gamma$ in the fibre. For a closed loop $\gamma$ (one with $\gamma(0)=\gamma(1)$), the composite transport around the loop is an automorphism of the fibre over the base-point: the \emph{holonomy} $\holo(\gamma)$. 

The phenomenon the main text calls ``return to the root, yet not to the origin'' is exactly a non-trivial holonomy: the loop closes in the base ($\gamma(0)=\gamma(1)$, one comes home to the same situation) while the transport in the fibre does not return the phase to its starting value ($\holo(\gamma)\neq \mathrm{id}$, one brings home a remainder). Where the structure group of the fibre is abelian---say $U(1)$, the natural choice if the phase of value is, as the name suggests, genuinely a \emph{phase}---the holonomy is a single angle, and by Stokes' theorem it equals the flux of the connection's curvature $\Omega$ through any surface $\Sigma$ bounded by $\gamma$ \parencite{berry1984,aharonov1987}:
\[
  \holo(\gamma) \;=\; \exp\!\oint_{\gamma} A \;=\; \exp\!\iint_{\Sigma}\Omega ,
\]
with $A$ the connection one-form and $\Omega = \mathrm{d}A$ its curvature. This is the precise content of two claims of the main text. First, that \emph{depth is the area a path encloses} (\xref{p09:sec:failures}): the accrued phase is the integral of the curvature over the enclosed surface, so a loop that encloses no area accrues no phase, and the ``shortcut'' that shrinks the loop toward a point drives the holonomy to zero. Second, that the good circle's phase is \emph{endogenous} (\xref{p09:sec:geom-phase}): the holonomy is fixed by the curvature \emph{internal} to the region the relation traverses, not by anything injected from outside---which is the geometric form of the distinction between the internal and the stolen phase. The sign of the holonomy---the spiral's positive accrual against the circle's zero and the vicious cycle's net loss---is the sign of the integrated curvature.

\subsection{The structured flow on a manifold (SFM)}
\label{p09:sec:app-sfm}

The contractual apparatus of \xref{p09:sec:contract} admits a complementary formalization. Where the geometric phase concerns the \emph{flow} on a given manifold, the contract concerns the \emph{shaping of the manifold} on which the flow runs. Model the free generativity as a flow on a manifold $M$; unconstrained, the flow's reachable set is all of $M$, including regions that lead to exhaustion or extraction. A contract is then a structuring of $M$ into $M'$ by three kinds of operation, corresponding to the three the main text names. It may \emph{excise} a closed region (the impassable singularities---the absolutely forbidden harms---removed from the reachable set, $M' = M \setminus U$). It may \emph{alter the boundary} (the exit conditions: $M'$ a manifold-with-boundary whose edge marks where the flow may licitly leave). And it may \emph{reshape the metric or the connectivity} (governing toward whom value may flow: a change in the topology or the metric structure that redistributes the geodesics). The constrained flow on $M'$ is the \emph{structured flow on a manifold} (SFM): free within $M'$, but no longer able to reach the excised or sealed regions.

The justice of the contract is then, as \xref{p09:sec:contract-justice} claims, a property of the \emph{geometry} of $M'$ and not of any path upon it: one asks whether the curvature and topology of $M'$ are such that every participant has a reachable region of positive holonomy, or whether some position in $M'$ is a structural sink---a region all of whose outgoing geodesics drain toward a single locus (structural siphoning), or a point from which no positive-holonomy loop is reachable (the structural fuel-position). And the corrective of \xref{p09:sec:legislation-is-flow}---that the shaping of the manifold is itself a flow---is the observation that the map $M \mapsto M'$ is not given once but is itself a time-dependent, generative process: $M$ is re-shaped continually, and this re-shaping (legislation) is a flow subject to the same criteria as the flow upon $M'$ (adjudication). There is no fixed terminal manifold; there is only the ongoing, generative re-legislation of $M_t$.

\subsection{The branching-process interface}
\label{p09:sec:app-branching}

The recursive regeneration of \xref{p09:sec:thesis} connects the geometric picture to a stochastic one, by way of the spectral radius of a multitype branching process---the formal machinery developed in concert with the Value-Foam programme. Model the cycle of interpretation as a branching process whose ``individuals'' are interpretive acts, each begetting further interpretive acts. Let $\lambda$ be the spectral radius (the Perron root) of the process's mean matrix. Then poeticity, in the sense of \xref{p09:sec:poem}, requires \emph{both} super-criticality \emph{and} a non-trivial phase increment per generation:
\begin{itemize}
\item $\lambda > 1$ with a non-trivial per-generation phase: the true poem---interpretations beget interpretations (the cycle endures), and each is a variation accruing fresh holonomy (the spiral);
\item $\lambda > 1$ with zero phase increment: the slogan, the meme, the jingle---reproduction without sublimation, a super-critical process each of whose generations merely copies the last (the circle that nonetheless propagates);
\item $\lambda < 1$: the dead poem---the interpretive process is sub-critical and dies out, no one interprets again;
\item the false poem (\xref{p09:sec:false-poem}): a process that advertises endurance while accruing zero phase---formally a degenerate case in which apparent branching carries no holonomy.
\end{itemize}
The conjugation of the two criteria---the spectral radius $\lambda$ governing \emph{whether} the cycle endures, the geometric phase governing \emph{whether} its endurance is a spiral or a circle---is the technical heart of the sequel, and is named here only to mark the road.

\subsection{The two debts, restated formally}
\label{p09:sec:app-debts}

The main text incurred two debts (\xref{p09:sec:caveats}); they are restated here in their formal shape, as open problems rather than settled results.

\emph{The first debt} is the demand that the fibre $F$ carry a genuine periodic, additive, gauge-invariant structure---an $S^1$, or a more general group---failing which $\holo$ is metaphor and not mathematics. The honest identification of this structure for ``the phase of value'' is the central task of the Value-Foam study: one must exhibit the bearer's circle-structure (in what sense is the sublimation of real value a \emph{phase}, with a well-defined notion of return and increment?), establish its additivity (do phases accrued along concatenated paths add?), and verify gauge-invariance (is the holonomy independent of the arbitrary choice of fibre-coordinate, as a genuine observable must be?). None of this is discharged here; it is incurred openly.

\emph{The second debt} is the demand that ``sublimation $=$ positive phase'' be argued and not stipulated. One must characterize, independently of the geometry, a quantity of sustainability---a \emph{margin of reproduction}: the net regeneration, per cycle, of the conditions of generation, in excess of what is consumed---and then \emph{prove} that, under the bundle formulation above, the sign of this margin corresponds to the sign of $\holo$. To weld the two by definition would be circular (a vicious circle in the most literal sense). The correspondence is, at present, a conjecture with strong motivation and no proof; the main text flags it as such (\xref{p09:sec:turn}), and the sequel must either establish it or revise the identification.

These two debts are not blemishes to be concealed but the precise form of the honesty the paper's epistemology requires (\xref{p09:sec:no-metric}): the geometric framework offers a structural \emph{characterization} of what the good circle is, and the formalization of that characterization is a programme in progress, not a finished edifice. To present it as finished would be to commit, in the appendix, the settling failure the whole paper is written against.
       % Appendix A: geometric phase, technical sketch

% ============================================================
\printbibliography[heading=bibintoc]

\end{document}
